\subsection{Reachability and repeated colors in joint source guards}
\label{app:core-injective-blockers}

Consider a bit query $(i,p)$, where $i$ is an input and $p$ is a proper
partial coloring of the other inputs. Write $p_r$ for its set of inputs
colored $r$. The query asks whether source $i$ contributes at $p$, in
the sense of \cref{core:thm:source-completion}. Let $Z$ be a finite set
of queries $(j,t)$ required to have bit zero. We call $j$ its head.
The heads, colors and inputs may recur freely across queries.

For $S\subseteq Z$ and each owner $r$, draw an edge $u\to j$ for
every $(j,t)\in S$ with $t_r=\{u\}$. Let $D_r(S;p)$ consist of the
vertices reachable, including by an empty path, from
\[
 p_r\ \cup\ \{j:(j,t)\in S,\ t_r=\varnothing\}.
\]
Thus the reachability tests use only the specified source bits.

\begin{theorem}[Reachability for joint source guards]
\label{core:thm:injective-blockers}
Let $K\ge3$. Suppose every negative query $(j,t)\in Z$ has
$|t_r|\le1$ for each owner $r$. A core word supplies $(i,p)$ and
omits every query in $Z$ if and only if every $S\subseteq Z$ satisfies
\begin{equation}\label{core:eq:injective-blockers}
 i\notin\{j:(j,t)\in S\}
 \quad\text{or}\quad
 \exists (j,t)\in S,\ r\in I:\ t_r\not\subseteq D_r(S;p).
\end{equation}
When this holds, one can start at the leaf $i$ and attach at most $|Z|$
siblings. Each attachment by $\star_r$ uses as its sibling a nonempty
$\star_r$-fold of inputs.
\end{theorem}
\begin{proof}
To read $(j,t)$, choose any owner $c$ missing from $t$. Give $j$ the
menu $I$, every colored input its prescribed singleton, and each other
input $I\setminus\{c\}$. Only $j$ supplies $c$, and the singleton
coalitions at this valuation are exactly the $t_r$, since $K\ge3$.
Thus membership of $c$ reads the source bit by the source formula
of \cref{core:thm:source-completion}. Call $c$ the mark of this row.
A subword supplying it cannot be a singleton: the only possible such
singleton would be $\{c\}$, and no input is that singleton.

Suppose \eqref{core:eq:injective-blockers} fails at $S$, and take a
shortest word witnessing the requirements for $S$. A leaf supplying the
positive must be $i$, which occurs as a head in $S$. At an $r$-root,
both children lose the $r$-singleton test on $p_r$, since the positive
mark survives. We claim that for each $v\in D_r(S;p)$ both children
also lose that test on $\{v\}$. This holds on $p_r$ by monotonicity.
If $t_r=\varnothing$, neither child can absorb at the corresponding
negative query. If $t_r=\{u\}$ and the claim holds at $u$, neither
can absorb there either. In both cases both children must omit that
row's mark. A child winning the $r$-test on $\{j\}$ would, by
singleton-profile absorption
(\cref{core:lem:singleton-profile-absorption}), either be the
$r$ singleton there or contain the entire menu of $j$, including its mark. Both alternatives
are excluded, proving the claim at the head $j$ and hence along
every path. Since $t_r\subseteq D_r(S;p)$ for every row of $S$ and
$|t_r|\le1$, neither child absorbs at any of these rows. Both omit
all their marks and one supplies the positive mark, contrary to
minimality. A named union likewise passes every omission to both children.

Conversely, induct on the number of remaining rows $S$, starting with
$S=Z$. If $i$ is absent from its heads, use the projection $i$.
Otherwise choose an owner $r$ from the second alternative and let
$T=V\setminus D_r(S;p)$. The $r$-fold of $T$ is nonempty and does
not absorb at the positive query. At every row with
$t_r\cap T\ne\varnothing$ it absorbs. At each remaining row its head
is in $D_r(S;p)$, by the initial-vertex or edge rule, so the fold
itself omits the mark. At least one row was absorbed. The induction
hypothesis gives a word for the remaining rows and the same positive.
Attaching the fold to this word by $\star_r$ enforces every removed
row. On the remaining rows neither operand supplies the mark, and at
the positive the inductive word supplies it and the fold cannot erase
it. This gives the stated spine.
\end{proof}

For a full family of antitone source predicates agreeing across observers,
take $Z$ to be all inclusion-minimal zeros at their respective heads.
If the family has a positive bit and every such zero uses each color
at most once, \eqref{core:eq:injective-blockers} on the maximal positive
bits is necessary and sufficient for realization. Take the union of
their witnessing spines. Antitonicity settles all other bits, and
source decomposition then settles every winner valuation. A family
with no positive bit cannot be realized: any leaf of a word transmits
its mark when all passive inputs are uncolored.

The test is polynomial in the explicitly listed positive and negative
queries. For each positive start with all negatives. Whenever the
second alternative applies, delete all rows on which the chosen fold absorbs.
Reachable sets can only shrink under deletion. Thus every subfamily
failing \eqref{core:eq:injective-blockers} stays present throughout:
none of its tails can meet the complement of the larger reachable set.
The process either exposes such a subfamily or reaches a set without
head $i$. It takes at most $|Z|$ rounds of ordinary graph reachability.

\begin{corollary}[Cuts on the input set]\label{core:cor:input-guard-cuts}
Under the hypotheses of \cref{core:thm:injective-blockers}, the
positive query is impossible exactly when some $H\subseteq V$ contains
$i$ and satisfies the following condition for every owner $r$ and every
nonempty $U\subseteq H\setminus p_r$:
\[
 \exists (j,t)\in Z:\quad j\in U,\qquad
   t_r\cap U=\varnothing,\qquad t_s\subseteq H\cup p_s\quad(s\in I).
\]
For $n=|V|\ge3$, every impossible query already fails with at most
$n\min(K,n)$ of the given negatives. The optimal worst-case bound
is $\Theta(n\min(K,n))$, uniformly in $K$ and $n$.
\end{corollary}
\begin{proof}
Let $S$ violate \eqref{core:eq:injective-blockers} and let $H$ be its
heads. Then $D_r(S;p)=p_r\cup H$: every head is reached either as an
initial vertex or from its reached tail. A first entry into $U$ is
therefore witnessed by a row of the displayed kind, including an
empty-tail row if it supplies the initial vertex.

Conversely, take all rows with $j\in H$ and $t_s\subseteq H\cup p_s$
for every $s$. If any vertices of $H$ are unreached at $r$, their set
$U$ satisfies the displayed cut condition. Its witness is either an
empty-tail row or an edge from a reached vertex, so reaches a member
of $U$, a contradiction. Every vertex of $H$ is a head: for $j\in H$
apply the condition to $U=\{j\}$ at an owner with $j\notin p_r$.
Thus these rows violate \eqref{core:eq:injective-blockers}.

Choose one row at each head $j$ of $S$. For every color $r$ present
in that row with $j\notin p_r$, also retain the row first reaching $j$
in the original $r$-search. For all other colors the chosen row makes
$j$ an initial vertex or $j\in p_r$. Induction in discovery order
therefore preserves every reachable set. Each chosen row has at most
$\min(K-1,n-1)$ colors, so at most $n\min(K,n)$ rows suffice.

For the lower bound take $d$ owners in a set $R\subsetneq I$, with
$2\le d<n$. Use inputs $x_1,\ldots,x_m$ and $u_s$ for $s\in R$,
where $m=n-d$, and query $x_m$ at the empty coloring. At each $x_j$
give one row for each $r\in R$: its $r$-tail is $x_{j-1}$ if $j>1$
and empty if $j=1$, and its $s$-tail is $u_s$ for every $s\in
R\setminus\{r\}$. The only row at $u_s$ has $s$-tail $x_m$.
All unspecified colors are empty. The full family violates the
reachability criterion: at each $r\in R$ the $r$-rows reach
$x_1,\ldots,x_m,u_r$ in order, and every other $u_s$ is initial.
In any obstruction $x_m$ must be reached at every $r$. Its first
arrival precedes any visit to $u_r$, whose sole incoming $r$-edge is
from $x_m$. The only route to $x_m$ avoiding $u_r$ starts with the
empty $r$-tail at $x_1$ and uses the $r$-row at every subsequent $x_j$.
Thus all $dm$ rows at these heads are required. Since $d\ge2$, at $x_1$
each $u_s$ occurs as an $s$-tail in a required row of another owner.
Reaching that tail forces the sole row at $u_s$. All $d(m+1)$
negatives are therefore necessary. Take
$d=\min(K-1,\max(2,\lfloor n/2\rfloor))$ to obtain the claimed order.
\end{proof}

\begin{corollary}[Complexity of the sample test]\label{core:cor:guard-cut-complexity}
For every fixed $K\ge3$, feasibility of one positive query and
explicitly listed negatives as in \cref{core:thm:injective-blockers}
is $\mathrm P$-complete under logspace reductions, even with at most two
colored inputs per negative and an acyclic graph of dependencies.
\end{corollary}
\begin{proof}
The deletion algorithm above gives polynomial time. Reduce monotone
circuit value \cite{Goldschlager1977} to infeasibility. Give each
circuit vertex a variable and take the output as the positive head with
$p$ empty. A true input has an empty-tail negative, a false input none.
An AND gate has one negative with its two predecessors colored $a,b$;
an OR gate has two alternatives, each with one predecessor colored $a$.
A gate with identical predecessors uses a single one-tail row.
The construction uses logarithmic workspace. Any obstructing $H$ has
at each head a row with all tails in $H$, by the singleton cut at that
head. Induction in circuit order makes every head in $H$ true.
Conversely, if the output is true, take all true vertices for $H$.
In any nonempty $U\subseteq H$, its earliest vertex has a row whose
tails are true and precede it, so the row meets the cut condition at
every owner. Thus the output is true exactly when the query is
infeasible. Closure of $\mathrm P$ under complement gives the stated
hardness for feasibility.
\end{proof}

\paragraph{When a color repeats.}
The same graphs eliminate all rows with singleton color classes even
when a witness must win specified coalitions. This bounds the remaining
search by the number of rows in which a color repeats. A single such row,
of arbitrary size, still admits a polynomial test.

A repeated color can already invalidate the criterion on three inputs.
For distinct $a,b$ consider $x\star_a(y\star_b z)$. Its minimal zeros are
\[
\begin{array}{c|l}
 x & y=z=a\\
 y & x=a\quad\text{or}\quad z=b\\
 z & x=a\quad\text{or}\quad y=b .
\end{array}
\]
At the uncolored positive for $x$, even replacing singleton edges by
the implications $t_r\subseteq D_r\Rightarrow j\in D_r$ gives
$D_r=\{x,y,z\}$ for every $r$ on this whole family. For $r=a$ the two
$b$ rows initially give $y,z$, and their joint implication gives $x$.
For $r=b$ the empty tails already give all three heads, as they do for
every other owner. The inner word wins the $a$-test on $\{y,z\}$
although it loses on both singleton coalitions. This is the step
in the necessity proof that changes when a color repeats.

Write $r:A$ for the demand $\Phi_r(A)=1$. We use three tests on a set
$S$ of negative source queries. Each asks for one word omitting all
their marks and satisfying one of the conditions below. The last column
gives a set of vertices $L_r$ for each prospective gate owner $r$:
\[
\begin{array}{c|l|l}
 &\text{condition}&L_r\\ \hline
 M &\text{the positive source query }(i,p)&p_r\\
 R &s:A_s\ (s\in J),\ \text{one coalition per owner}
     &\displaystyle\bigcup_{s\in J\setminus\{r\}}(V\setminus A_s)\\
 G &k:A,\ \Phi_k(B)=0
     &B\ \text{if }r=k,\quad V\setminus A\ \text{otherwise}.
\end{array}
\]
The empty selection in $R$ asks for any bounded word. Write $h(S)$
for the heads occurring in $S$. Let $D_r$ be
the vertices reachable on the edges $u\to j$ with $t_r=\{u\}$,
starting from $L_r$ and all heads $j$ of rows $(j,t)\in S$ with
$|t_r|\ne1$. Put $S_r=\{(j,t)\in S:t_r\subseteq D_r\}$.

To describe the test, we also need to express conjunctions of demands
in terms of these three forms. Given a finite family $P$ of demands
$s:A$, require $R$ for every subfamily with pairwise distinct owners,
including the empty subfamily. If a mark $(i,p)$ is also requested,
require $M$ for that mark and $G$ on $(s,A,p_s)$ for every $s:A$ in
$P$. If a losing coalition $\Phi_k(B)=0$ is requested instead, require
$G$ on $(k,A,B)$ for $A=V$ and for each $k:A$ in $P$.
All tests use the same negative rows.

\begin{theorem}[Recognition with repeated color classes]
\label{core:thm:repeated-blocker-budgets}
For each of $M,R,G$, the following rules decide feasibility.
\begin{enumerate}
\item A projection satisfying the rows and the requested condition succeeds.
\item If $S_r\subsetneq S$ for some $r$, replace $S$ by $S_r$,
keeping the condition. Any such $r$ may be used.
\item Otherwise try each distinct $(r,A)$ with $A=t_r$ and $|A|\ge2$
in some row of $S$ such that $A\setminus L_r$ meets $h(S)$.
Drop all rows with $A\subseteq t_r$ and add the
demand $r:A$ to the requested condition. Test the resulting conjunction
as above. Success requires at least one successful choice.
\end{enumerate}
Only repeated classes containing a negative head can branch, and each
such pair $(r,A)$ or row can be used at most once along a branch.
In particular, arbitrary repetitions among inputs that are never heads,
together with arbitrary singleton dependencies between heads, admit a
polynomial test.
\end{theorem}
\begin{proof}
First justify the conjunction tests. Take one representative of every
term operation omitting all specified marks, giving a finite set $H$.
It is closed under owned choice up to equality of operations. For each
$s:A$ let $E_{s,A}=\{h\in H:\Phi_s(h)(A)=1\}$. The least prospective
profile at owner $s$ on the inputs $H$ has as winning coalitions
the upsets generated by its required $E_{s,A}$ and by $H$ itself.
The rainbow tests say exactly that these profiles have the common-source
property of \cref{core:thm:profiles}, including $H\ne\varnothing$.
They therefore synthesize a term in these bounded inputs meeting all
demands. To make $B$ losing at $k$, every generator of the required
$k$-upset must escape $E_{k,B}$, precisely the stated gap tests.

For the mark query use the valuation in the proof of
\cref{core:thm:injective-blockers}, with marked winner $c$.
Test $M$ puts $c$ in the union of the bounded inputs and ensures that $H$
escapes every $E_{s,p_s}$: a value containing $c$ cannot be a singleton.
The remaining gap tests make
each required $E_{s,A}$ escape $E_{s,p_s}$ too. The synthesized term
therefore has no atomic output on the query. Its greatest profile
completion (\cref{core:prop:profile-maximum}), still composed of
bounded inputs, supplies their union there. Conversely, a common
witness supplies all the intersection, gap and mark tests.

For the second rule put $T=V\setminus D_r$. Attaching its $r$-fold
to a witness for $S_r$ absorbs on the deleted rows; each remaining
head lies outside $T$, so the fold omits its mark. For $M$, the fold
cannot absorb at $p$. For $R$, $T\subseteq A_s$ for each foreign
demand, and the witness retains the own-owner demand at the attachment.
For $G$ at $r=k$ the fold loses on $B$, while at $r\ne k$ it wins
the foreign test on $A$. Thus attachment preserves the condition in
each case. Restriction proves the converse.

At the third rule a successful branch supplies an atom $\{r\}$ on every
dropped row. Conversely, every current $r$-class lies in
$D_r\subseteq L_r\cup h(S)$, since every reachability edge ends at a head.
If a class $A$ lies in $L_r$, no witness can win $r:A$. For $M$ it
would absorb at $p$. For $R$ the class misses the intersection of the
foreign demands, contrary to the common-source property. For $G$ it
is contained in the losing coalition $B$ if $r=k$, and is disjoint
from the winning coalition at $k$ otherwise. Thus every winning repeated
class passes the branching condition.

Suppose now that there is a shortest witness but no branch succeeds.
No witness can win on a repeated $t_r$, since it would pass that branch.
At an $r$-root both children
therefore fail to absorb on the $r$-repeated rows, and also on rows
with empty $t_r$. They omit those marks. By singleton-profile
absorption, both lose the $r$-profile on each of these singleton heads.
Both also lose on every singleton in $L_r$: for $M$ this follows from
survival of the mark, for a foreign demand from the common-source
property, and for the own-owner gap from loss on $B$. Propagating along
the edges as in \cref{core:thm:injective-blockers} proves singleton
loss throughout $D_r$. As $S_r=S$, neither child absorbs on any of
the remaining singleton-tail rows. Both children consequently meet
all the negative rows. One inherits the primitive condition: a surviving
mark comes from a child, foreign positive demands pass to both, and for
a same-owner gap the OR/AND profile rules preserve both signs in one
child. This contradicts shortestness unless a projection succeeds.
Each branch in the third rule removes a row containing the chosen
repetition and all occurrences of that pair.
\end{proof}

On an explicitly listed sample the first two rules take polynomial
time between branches: each strict restriction drops at least one row.
If only one row has a repeated class meeting $h(S)$, at most $K$
choices need be tried and none of the smaller budgets branches.
More generally, let $d$ count the distinct repeated pairs $(r,A)$
with $A\cap h(S)\ne\varnothing$ in the initial sample. The running time
for a marked query is $2^{O(d\log(d+1))}$ times a polynomial in its
explicit length. Indeed each choice needs only two smaller tests:
the original condition and that condition with its $r$-entry replaced
by $r:A$, in the assignment notation of \cref{core:thm:pointed-guard-cut}
below. Their $r$-choice meets the original condition and the added
demand; conversely a common witness solves both tests. At most $d$
choices branch at each of at most $d$ levels.
For a marked query $d$ can be reduced further. Its projection
suffices if $i\notin h(S)$. Otherwise apply
\cref{core:lem:unmarked-background} with $H=h(S)$ first:
keep just the rows whose postings outside $H$ are contained in $p$,
and erase those outside postings. It now suffices to count distinct
$(r,t_r\cap H)$ of size $|t_r\cap H|\ge2$. In particular, each
color may occur arbitrarily often in a row without any branching
if at most one of its postings there is at a negative head. Testing
each maximal positive bit against the minimal zeros decides an
explicit joint guard table.

\paragraph{Obstructions on subfamilies of the given queries.}
For arbitrary repetitions one can use just one coalition per owner.
Fix the positive $(i,p)$ and a set $Z$ of negatives. A partial assignment $a$ selects,
for each owner in its domain, a nonempty coalition occurring at that
owner in $Z$. Define the eligible inputs by
\[
 L(a)=
 \begin{cases}
 \{i\},&\operatorname{dom}a=\varnothing,\\
 a_r\setminus p_r,&\operatorname{dom}a=\{r\},\\
 \displaystyle\bigcap_{r\in\operatorname{dom}a}a_r,
      &|\operatorname{dom}a|\ge2 .
 \end{cases}
\]
For a nonempty $r$-coalition $A$ of a row in $S$, put
$S_{r,A}=\{(j,t)\in S:A\not\subseteq t_r\}$. Write $a[r:=A]$
for insertion or replacement at $r$.

\begin{theorem}[A cut on subfamilies of source queries]
\label{core:thm:pointed-guard-cut}
The positive cannot be supplied while omitting $Z$ exactly when
there is a collection of pairs $(S,a)$ containing a pair with empty
assignment and satisfying the following conditions at every member:
\begin{enumerate}
\item $L(a)\subseteq h(S)$;
\item for each nonempty $r$-coalition $A$ occurring in $S$, the
collection contains a pair $(T,a[r:=A])$ with $T\subseteq S_{r,A}$.
\end{enumerate}
For the sample decision problem let $m=|Z|$, and let $W$ be the sum, over
owners, of the widths of their inclusion orders on the distinct
nonempty coalitions occurring in $Z$. Feasibility is decidable in
polynomial space, and in time $(m+1)^{2W}$ times a polynomial in the
explicit input size. In particular, at fixed $K$, arbitrarily many
rows and arbitrarily many repeated coalitions at every owner are allowed
in polynomial time if each owner's coalitions form a chain.
\end{theorem}
\begin{proof}
With an empty assignment ask for the mark. With a single entry $r$
ask for a word winning the $r$-test on $a_r$ and losing it on $p_r$.
With two or more entries ask for all their positive singleton tests.
A projection satisfies the assignment and omits the rows of $S$
precisely when its input is in $L(a)\setminus h(S)$.

A nonempty coalition $A$ of owner $r$ gives the following construction.
On $S_{r,A}$ take a word for $a$ and a word for $a[r:=A]$.
Combine them by $\star_r$. This absorbs every removed row. For an
empty $a$, the second word loses on $p_r$ and the first transmits
the mark. If $a$ has the sole entry $r$, both words lose on $p_r$
and the first wins on $a_r$. If $a$ has a sole entry $s\ne r$, both
words win the $s$-test on $a_s$ and the first loses on $p_s$.
For multiple entries, all foreign tests pass to both operands and
the first supplies any required $r$-test. Thus the result meets $a$
on $S$ in all cases.

Conversely, take a shortest word meeting $a$ on $S$. If it is no
projection, it must yield a singleton on some row of $S$. Otherwise
the root is a union on every row, so both children omit all their
marks. One child retains the assignment: a mark comes from one child,
positive foreign tests pass to both, and the owned test is disjunctive.
For a gap the losing test passes to both children at its own gate and
to at least one at a foreign gate. Descent contradicts shortestness.
If the word yields $\{r\}$ on a row with $t_r=A$, then $A\ne\varnothing$,
its $r$-test on $A$ is positive, and the word itself meets both
assignments on $S_{r,A}$. When $a$ was empty this uses that a surviving
mark precludes an $r$-singleton at $p$. We have proved equivalence to
a fitting projection or a choice of $(r,A)$ with both smaller tests
feasible.

For each assignment separately, take all inclusion-minimal infeasible
subfamilies of $Z$. At any such $S$ the old assignment is feasible on
the proper subset $S_{r,A}$. The construction forces the updated
assignment to be infeasible there, giving the second cut condition.
If $Z$ fails the positive query this collection includes a member with
empty assignment. In the other direction, a feasible member of a cut
with the fewest rows cannot use a projection. Its singleton row gives
an updated assignment feasible on $S_{r,A}$ and hence on the smaller
member $T$, a contradiction.

The two smaller tests can be run in depth first order with at most $m$
strict row deletions. For the time bound, cache each pair of a remaining
budget and an assignment. At an owner $r$ the coalitions used for
deletion generate an upset in the poset of the original coalitions.
Its minimal elements form an antichain of size at most the width $w_r$,
so there are at most $(m+1)^{w_r}$ choices. These upsets determine the
remaining rows. Each active owner also has at most $m+1$ possible
assignment entries, counting absence, and there are at most $W$ active
owners. Each pair takes polynomial work once its children are known.
\end{proof}

\paragraph{Monochromatic cycles in the posting graph.}
Draw an edge $j\xrightarrow{r}u$ whenever a negative row at $j$ posts
$u$ at $r$. A cycle is monochromatic if all its edges have the same
color. Individual rows may post several colors.

\begin{lemma}[Backward induction on the input names]
\label{core:lem:ranked-posting-survivor}
Choose one row $t^j$ at each input $j$. If their posting graph has no
monochromatic cycle, every word supplies at least one head's mark on
its chosen row. Conversely, the fold on any monochromatic cycle, taken
at its color, omits every chosen mark. Whenever there is an omitting
word, its minimum possible number of leaves equals the length of a
shortest monochromatic cycle.
\end{lemma}
\begin{proof}
Give each owner $r$ a rank on the inputs, strictly decreasing along
the $r$-edges. Evaluate a word to an input name by choosing at each
$r$-gate the child whose resulting input has smaller $r$-rank,
breaking ties arbitrarily. If a subtree wins the $r$-test on a set
$E$, its resulting input has $r$-rank at most the largest rank in
$E$. This follows by induction: an $r$-gate can choose at least as
well as a winning child, and at every other gate both children win.

Let $j$ be the resulting input at the root. Follow the selected path
to its leaf; every selected child on this path also has result $j$.
An $r$-sibling cannot win on $t^j_r$: all its postings have rank
strictly below $j$, whereas the gate chose $j$. On a probe with full
menu at $j$ no sibling on the path can therefore absorb at its gate.
The leaf supplies the mark. This argument also works on an infinite
input set: restrict the rows and the ranks to the finitely many
names in the word. Conversely, in a cycle fold every used
head has a cycle successor posted at the fold owner and every other
head is absent. Restricting the obstruction to the support of an
omitting word gives a cycle within its leaves; a shortest cycle uses
each of its vertices once.
\end{proof}

For comparison, if all losses at a head use a single designated owner
$c(j)$, there is a direct extraction even with several losses there.
Assign any owner to inputs without losses.

\begin{lemma}[A single-owner fold in every feasible support]
\label{core:lem:monochromatic-safe-fold}
The support of every word omitting these losses contains a nonempty
set $S$ of inputs of one designated owner $r$ whose $\star_r$-fold
also omits them. In particular every inclusion-minimal feasible
support has a read-once realizer of this form.
\end{lemma}
\begin{proof}
Write $H_r=\{j:c(j)=r\}$. Some $r$ has
$\Phi_r(V\setminus H_r)=0$, since otherwise the positive coalitions
$V\setminus H_r$ would have empty intersection, contrary to
\cref{core:thm:profiles}. Keep both children of each $r$-gate and
one child losing this test at every foreign gate. The resulting
$r$-fold has all its leaves in $H_r$. Test a pure $r$-row with the
full menu at its head: the source bit is independent of the other
winners in that menu by source decomposition. Every foreign gate
then acts by union, and $\star_r$ is monotone, so this pruning
preserves omission there. On a foreign row its head is
absent from the fold. Idempotence removes repeated inputs.
\end{proof}

\begin{theorem}[A cut with one loss per head]
\label{core:thm:monochromatic-head-cut}
Suppose every head with several losses belongs to monochromatic cycles
of at most one color. Let $(i,p)$ be a positive query for which no
vertex of an $r$-cycle is posted at $r$ in $p$. These hypotheses
allow arbitrary mixed rows. The first is automatic when each head
has one loss, or all its losses use only its designated owner.
The query is impossible exactly when one
can choose a set $H\ni i$ and one loss $(j,t^j)$ at each $j\in H$
such that
\[
 t^j_r\setminus H\subseteq p_r\quad(j\in H,\ r\in I),\qquad
 \text{the chosen rows on $H$ have no monochromatic cycle}.
\]
The query is decidable in polynomial time from either the tails or
the canonical prime guards. A possible query has a witness with at
most $1+K|V|^2$ leaves.
\end{theorem}
\begin{proof}
For the obstruction direction restrict to the chosen rows.
By \cref{core:lem:unmarked-background}, a witness would give a word
on $H$ omitting every restricted loss, contrary to
\cref{core:lem:ranked-posting-survivor}. This direction requires
neither hypothesis.

For the converse, on the remaining rows take the greatest support
$S_r\subseteq V\setminus p_r$ of a safe $r$-fold, allowing empty
support. Start with $V\setminus p_r$ and remove an input whenever
one of its rows has no $r$-posting in the current set. A fold must
absorb whenever it uses the marked head, so this computes $S_r$.
Delete the rows with an $r$-posting in $S_r$. The fold omits all
remaining marks and cannot absorb at the positive. Attaching it by
$\star_r$ therefore lifts any witness after deletion to one before
deletion. Repeat over the owners. If $i$ no longer heads a row,
its projection succeeds.

At stabilization let $H$ be the remaining heads, with $i\in H$.
Every tail posting outside $H$ agrees with $p$, since its unheaded
projection could otherwise have deleted the row. For each $r$ start
with the heads on $r$-cycles in the remaining graph restricted to $H$.
Iteratively remove any head having a remaining row whose $r$-postings
miss the current set. If a nonempty set remains, choose one $r$-posting
in that set from every row at each surviving head. A sink strongly
connected component has an edge out of each vertex. Every vertex
therefore lies on a monochromatic cycle and is outside $p_r$.
The component meets every tail at its heads, so its fold is safe
and hits a remaining row. This contradicts stabilization. Thus every
head in the starting set is removed. At each such head choose its
removing row. The choices for different colors agree wherever they
overlap, since a head lying on cycles of two colors has just one row.
Choose any remaining row at every other head. An $r$-cycle would
consist entirely of heads from the $r$-search, and each edge on it
leads to an earlier removal. This proves the required acyclicity.

To implement deletion, record $U_r$, the union of the fold supports
used at $r$. A row remains when its $s$-postings miss $U_s$ for
every $s$. In the support test, further exclude the $r$-postings
on the current candidate $S_r$. A head can remain in $S_r$ exactly
when no zero survives these exclusions. Scan the listed rows, or
on each observer facet test whether the allowed box is positive by
containment of a listed prime in its set of exclusions. Each strict
increase of a register adds an input-owner pair. At most $K|V|$
folds, each with at most $|V|$ leaves, and a final projection suffice.
\end{proof}

The same criterion handles prescribed singleton tests when the
augmented graph satisfies its hypotheses. For a required win $s:A$,
add the pure $s$-row posting $A$ at each head outside $A$, as in
\eqref{core:eq:total-monochromatic-source}. For a required loss
$\Phi_r(B)=0$, some input outside $B$ must supply a mark when just
$B$ is posted at $r$: give every unposted input the full menu and
apply source decomposition. Thus one may test each required mark
and each losing test separately on the augmented rows, and take
the union of their witnesses. This preserves every required singleton
win and every omission. If there are only winning tests, any feasible
word has a source at the empty coloring.

With just one loss at each headed input this is a reachability test.
From $i$ follow the edges $j\xrightarrow{r}u$ with $u\notin p_r$.
The mark is possible exactly when this search reaches an unheaded
input or a monochromatic directed cycle.
Indeed every cut must include the entire reachable set. Conversely
that set itself is a cut if neither event occurs: every edge of a
monochromatic cycle inside it would have been traversed, by the
hypothesis on $p$. In particular the uncolored test allows arbitrary directed
cycles and arbitrarily large tails.

One posting at its own owner can already change this answer. Give
$x,y,z$ respective single posting colors $a,b,b$ and tails
$\{z\},\{x,z\},\{y\}$.
Query $y$ with only $z$ posted at $b$. The search reaches all heads
but no monochromatic cycle, since it skips $y\to z$. Nevertheless
$y\star_b(z\star_a x)$ supplies the mark. The inner gate absorbs
on the $x$-row, both inner inputs are $b$ on the $y$-row, and $y$
is $b$ on the $z$-row.

The restriction on vertices belonging to cycles of several colors is
also sharp. Take the two reciprocal pairs $x\leftrightarrow_a y$ and
$x\leftrightarrow_b z$, with a separate one-post loss for each edge.
Choosing either row at $x$ leaves a monochromatic two-cycle.
Nevertheless no word omits all four losses, by
\cref{core:thm:additive-greatest-minorant}. At an uncolored query
no closed cut can avoid both pairs: every remaining head points to
$x$, or already is $x$, and whichever row is chosen at $x$ returns
to it in the same color.

When there is no monochromatic directed cycle anywhere, the cut
also tests a \emph{complete} guard family without quantifying over
positive colorings. Choose a nonempty set $H$ of heads and one row
$t^j$ at each. Put $B_r=\bigcup_{j\in H}(t^j_r\setminus H)$.
For a family with some positive bit, realizability is equivalent to
the following condition on every such choice: if the $B_r$ are pairwise disjoint and
some $B_r$ is empty, then for every $j\in H$ some loss row at $j$
is contained in $B$ componentwise. Indeed the $B_r$ give the least proper
coloring compatible with that boundary. If a head is positive there,
the cut forbids its mark. Conversely any impossible prescribed
positive gives a cut, and remains positive on restricting its
coloring to this boundary. If every bit is zero there is no term;
otherwise union of the individual witnesses realizes the family.

\begin{corollary}[Recognition by closing the endpoints of an edge]
\label{core:cor:monochromatic-boundary-recognition}
Suppose each headed input has one loss, the graph has no
monochromatic directed cycle, and some bit is positive. For every
edge $j\to u$ and every owner $k$, close $\{j,u\}$ under two rules.
Include every vertex posted at $k$ by an included head. Include a vertex
whenever included heads post it at two distinct colors. The
family is realizable exactly when every such closure reaches an
unheaded input. In particular full recognition takes polynomial time.
\end{corollary}
\begin{proof}
A proper boundary cannot enclose an edge: if a row at $j$ posts in $H$, its
head is positive at the boundary, and otherwise it is suppressed
there. For a chosen missing owner $k$, $B_k=\varnothing$ is the
first closure rule, and disjointness of the $B_r$ is the second.
Hence the least closed set containing the edge is a forbidden set
of heads precisely when it avoids the unheaded inputs. There is one
finite closure computation per edge and owner.
\end{proof}

Specialize now to heads with a designated owner $c(j)$: every loss
at $j$ posts its tail $T_j$ entirely at $c(j)$.
With at most two owners among the heads and no same-owner edge,
there is also a polynomial test on explicitly listed minimal tails,
with no bound on their number per head. For each tail $T_j$, each
headed $u\in T_j$, and each tail $T_u$, require $T_j\cap T_u\ne
\varnothing$. Any shared input is unheaded: a headed one would
receive an edge from its own owner. Thus a violating cut cannot have
an internal chosen edge $j\to u$ and satisfy the pair condition,
since its two tails would conflict outside. Conversely disjoint
tails give the cut $\{j,u\}$ with a proper boundary. Head $j$ is
positive there since its postings at $c(j)$ are just
$T_j\setminus\{u\}$, which contains no minimal tail at $j$.
Throughout, a family with all bits zero is rejected separately.

On prime lists alone, this two-owner pair test is polynomial at
width two and coNP-complete at width three. To find disjoint tails
with $u\in T_j$, choose a prime $G$ at $j$ containing $u$; all
primes here are viewed as sets of inputs on a facet seeing the
head's owner. Seek a set $X$ meeting every prime at $j$, with
$X\cap G=\{u\}$, whose complement meets every prime at $u$.
Minimizing the two hitting sets yields disjoint minimal tails and
preserves $u$ in the first. Conversely minimality of $T_j$ gives
such a $G$. Each prime gives one clause of its own width in the
membership bits of $X$; the prescribed intersection only adds unit
clauses. Thus width two gives a polynomial number of 2SAT tests.

For hardness, use just two heads $j,u$ at different owners, and
two unheaded literal inputs for every variable of a 3CNF formula.
At $j$ use the primes $\{u\}$, the literal pair of each variable,
and the sets of literals in each clause. At $u$ use just the
literal pairs. Remove redundant supersets to make the lists
canonical, and give the facets hiding the respective owners the
empty guard. The sole conditional dependency is $j\to u$. The
first hitting set and its disjoint counterpart choose opposite
literals of each variable; the first must meet every clause.
Hence an obstructing pair exists exactly for a satisfiable formula.

\begin{theorem}[Two losses per head suffice for coNP-completeness]
\label{core:thm:two-monochromatic-losses}
At each fixed $K\ge3$, full recognition is coNP-complete even with at
most two monochromatic losses per head, no edge between heads of the
same owner, and at most two exclusions per prime. Both the minimal
losses and the canonical primes may be supplied.
\end{theorem}
\begin{proof}
The boundary criterion gives a certificate of failure: list the heads,
one loss per head and a head positive at the resulting proper boundary.
With prime lists alone, a proposed minimal loss can be verified by
testing it and its one-post deletions. The family with no positive bit
is recognized directly.

Reduce 3SAT to failure of recognition. Use three owners $a,b,c$.
Give each Boolean variable two unheaded inputs for its two literals.
There are three control heads, one per owner. Call a further head
mandatory if the construction below requires every violating cut to
contain it. Each control has one tail containing the other two
controls and all mandatory heads of a different owner. The $b$-control
also posts all optional heads. Every tail at a noncontrol head contains
the two controls of owners different from its own. In addition it
posts the one input specified by the following choices.

For each variable introduce two mandatory heads, one at $a$ and one
at $b$. Each chooses either literal input of the variable.
For each clause with literals $\ell_1,\ell_2,\ell_3$, introduce a
mandatory $a$-head whose two choices are the input for $\ell_1$ and
an optional $c$-head $z$. The two choices at $z$ are optional
$a$-heads $w_2,w_3$. The sole choice at $w_j$ is the input for
$\ell_j$. Here a choice denotes one possible tail in the loss list.
All tails are monochromatic and no edge joins heads of the same
owner. On a facet seeing the loss color, a prime excludes a common
posting or excludes the private postings of both alternative tails.
A head with one tail has only single-literal primes; on the facet
hiding the color the empty guard suffices. Thus both presentations
have polynomial size.

A violating cut contains an edge of one of its chosen tails: otherwise
those tails alone suppress every selected head at the boundary.
The endpoints have distinct owners. Both post the third control, which
must therefore lie inside the cut. These three heads now force every
control inside: the two heads of owners different from a missing
control both post it. The controls in turn force every mandatory head,
since two controls of different owners post it. The opposite-color
choices at each variable pair must choose different literal inputs.
Declare the $a$-choice true. In a clause the first choice either
posts a true literal at $a$ or selects $z$. In the latter case $z$
must lie inside: outside it would also be posted at $b$ by the
control. The same argument puts the chosen $w_j$ inside, and its
$a$-posting makes $\ell_j$ true.

Conversely, given a satisfying assignment take all controls and
mandatory heads, choosing according to the assignment. In each clause
use a true literal, adjoining $z$ and its selected $w_j$ precisely
when using one of the last two literals. The boundary posts omitted
optional heads at $b$, true literal inputs at $a$, and false ones
at $b$. It is consistent and omits $c$. All controls are positive
there since their tails contain other controls. This gives the
violating cut.
\end{proof}

\begin{proposition}[A single posting missed by every boundary cut]
\label{core:prop:monochromatic-cut-obstruction}
For every $K\ge3$ a seven-input family with one monochromatic loss per
head passes all acyclic boundary cuts, even when every owner may occur
on the boundary. It has idempotent pure-owner closures and admissible
singleton profiles. Nevertheless its losses force a singleton when
just one input is posted, where the family prescribes another winner.
Only one pair of inputs forms a monochromatic cycle.
\end{proposition}
\begin{proof}
Use distinct owners $d,\alpha,\beta$. In the table a string denotes
a set of input names; upper- and lower-case names are distinct.
\[
\begin{array}{c|ccccccc}
j&A&a&X&B&Y&b&Z\\ \hline
c(j)&\alpha&\alpha&d&\beta&d&\beta&d\\
T_j&XBbZ&BbZ&ABY&AaYZ&aXB&AXZ&Ab
\end{array}
\]
The two $d$-heads $X,Y$ post one another, so their fold lies below
the family. There is no other same-owner edge. Each pure closure
adds its heads in one step: adding $X$ requires $Y$ already present,
and conversely; $Z$ is posted by neither. The $d$-profile requires
$A,a,B,b$ and at least one of $X,Y$. Each foreign profile requires
all inputs except the heads of its own owner; unused owners require
all inputs. Thus every positive rainbow has a common input in $XY$.

A consistent boundary must include every common posting of two
included heads of different owners. The following implications
list the unordered edges between different owners. Each pair on
the left forces the indicated pair by this one rule, and the rows
are ordered so that each implication leads back to $XY$.
\[
\begin{array}{c|c@{\qquad}c|c}
 XB&XY & AX&XB\\
 Ab,\ Xb&AX & AZ,\ bZ&Ab\\
 AB,\ BZ&AZ & ab,\ aZ&bZ\\
 aB&aZ & aY,\ BY&aB
\end{array}
\]
Consequently a set enclosing any edge either contains the two-cycle
$X,Y$ or has inconsistent boundary. This verifies every acyclic cut.

Write $P=\{X\}$. Omission of all $s$-rows at a set $L$ of $s$-heads
whose tails avoid $L$ forces $\Phi_s(V\setminus L)=1$: post everything
outside $L$ at $s$. Each head in $L$ is suppressed even with its
menu free, so only the singleton $s$ can result. This applies to
$Aa$ at $\alpha$, $Bb$ at $\beta$, and to $XZ,YZ,Z$ at $d$.
It rules out absorption on all four foreign rows. For example, an
$\alpha$-win on $T_A=XBbZ$ would have to meet simultaneously the
forced $d$-coalition excluding $XZ$ and the $\beta$-coalition
excluding $Bb$. The same argument covers $T_a,T_B,T_b$, using their
indicated foreign and $d$ parts.

It remains to exclude absorption on the three $d$-rows under the
assumption $\Phi_d(P)=0$. We give the smaller sets of rows needed
in each step. Whenever absorption on all retained rows is excluded,
a gate acts by union there, so both children keep the omissions.
One child keeps a same-owner winning and losing pair of tests:
at that owner's gate both lose and one wins, and at a foreign gate
both win and one loses. Likewise one child retains wins at distinct
owners. Following these choices must end at a projection.
\begin{enumerate}
\item Even with just the rows at $A,b$, no word wins the $d$-test on
$T_Z=Ab$. Its intersections with $T_A,T_b$ are respectively $b,A$,
excluded by the opposite row's forced profile. Descent would select
one of the forbidden heads $A,b$.
\item With just the rows at $a,B$, a word losing on $P$ cannot win
on $T_Y=aXB$. Its intersections with $T_a,T_B$ are $B,a$,
respectively. The opposite row again excludes absorption, and the
only projections in $T_Y\setminus P$ are $a,B$.
\item Retain the rows $A,a,B,Y$. A $d$-win on $T_X=ABY$ excludes
absorption on $A,a$, since their tails meet $T_X$ only at $B$.
The second step excludes absorption on $Y$. If $B$ were atomic,
the word would win on both $T_B$ at $\beta$ and $T_X$ at $d$.
Retain only $A,a,Y$. The two winning coalitions intersect in $AY$,
disjoint from both $\alpha$-tails, and $T_B\cap T_Y=\{a\}$, excluded
by the $\alpha$-row at $a$. Descent with the two wins would end at
$A$ or $Y$, both forbidden. Thus $B$ cannot absorb either. Descent
with the original $d$-win and loss would end at $A,B$ or $Y$.
\end{enumerate}
On all seven rows, therefore, a word losing on $P$ could be atomic
nowhere. Descending with that loss would give a bounded projection,
although every input heads a loss. Hence every bounded word gives
the $d$-singleton when just $X$ is posted at $d$. The family
prescribes the mark of $A$ there. By \cref{core:lem:singleton-profile-absorption},
at any valuation each bounded word must either return $\{d\}$ or
include the entire menu at $X$.
\end{proof}

\begin{corollary}[The activation principal from the surviving inputs]
\label{core:cor:monochromatic-principal-box}
Let $D$ be the greatest word omitting the prescribed losses, or $0$ if
there is none. Assume the two cycle hypotheses of
\cref{core:thm:monochromatic-head-cut} for a proper coloring $p$.
Let $A$ consist of the inputs $j$ whose mark is possible at
$p|_{V\setminus\{j\}}$. For every proper coloring $b$,
\[
             b\preceq_D p\quad\Longleftrightarrow\quad
                    b|_A\subseteq p|_A .
\]
If $A\ne\varnothing$, a single word omitting the losses has support
exactly $A$ and supplies all these marks. Thus the tests of
\cref{core:thm:monochromatic-head-cut} compute the whole principal.
\end{corollary}
\begin{proof}
Grow a set $J$, initially empty. Discard losses headed in $J$ and
those whose postings on $J$ are outside $p|_J$; strip $J$ from the
remaining rows. If a safe $r$-fold on a nonempty set
$T\subseteq V\setminus(J\cup p_r)$ exists for this reduced sample,
adjoin $T$ to $J$. On its own support the fold supplies every mark
at $p$ and its principal there is ordinary subcoloring:
in \cref{core:thm:one-owner-activation-order} its closure of the
empty set is empty.

Maintain a word $B$ on exactly $J$ with these two properties,
omitting the original losses. The first fold serves as $B$. For a
later fold $F$, use \cref{core:thm:block-source-extension} over $B$
with coefficient $F$ on backgrounds $q\preceq_B p|_J$ and zero
elsewhere. This extension preserves the old source bits. A loss at
a new head either lies outside that principal or is omitted by $F$
on its stripped row. All other heads outside $J\cup T$ are absent.
At $p$ each old and new source survives, giving exactly the claimed
support. The greatest coefficient at $p|_J$ contains $F$.
By \cref{core:thm:block-source-extension,core:thm:activation-faithfulness}
the new principal requires subcoloring on both $J$ and $T$;
antitonicity supplies the converse.

At termination there is no nonzero safe word on the reduced sample.
An unheaded variable would give an admissible projection; otherwise
the removal argument of \cref{core:thm:monochromatic-head-cut}
applies on the reduced graph. It gives either a safe fold on an
$r$-cycle component avoiding $p_r$, or one row at each remaining
head with no monochromatic cycle. The former would have extended
$J$, and the latter excludes every word by \cref{core:lem:ranked-posting-survivor}.
Write $R=p|_J$. By \cref{core:lem:unmarked-background}, applied after
discarding the $J$-headed losses, no word omitting the original losses
can supply any outside mark even at $R$. Thus $J=A$. If $D=0$ the
principal is constantly true. Otherwise $D$ captures all outside inputs
under $R$ by \cref{core:cor:colored-cofactor-dependence}. Every interval
predicate above $D$ positive at $p$ is positive at $R$, so
\cref{core:lem:captured-rectangle} includes the entire displayed box.
The reverse inclusion follows from $B\subseteq D$ by
\cref{core:thm:activation-faithfulness}, extending the predicates of
$B$ over the unused inputs; for $J=\varnothing$ it is automatic.
\end{proof}


\paragraph{Independent color classes.}
Join the head of each negative row to every input colored in that row,
forming an undirected graph on $V$. On a triangle-free graph, each
neighborhood is independent, so every color class of every row is
independent. Only the latter property will be needed. For each
$S\subseteq Z$ and owner $r$, let $D_r$ be the least set closed under
the implications $t_r\subseteq D_r\Rightarrow j\in D_r$ for
$(j,t)\in S$, starting from the empty set. Put $T_r(S)=V\setminus D_r$.
This is the greatest support of a safe $r$-fold, allowing empty support:
a nonempty fold omits all marks exactly when every row whose head it
contains has an $r$-posting in its support.

\begin{theorem}[Winning an independent coalition]
\label{core:thm:triangle-free-neighborhood-profile}
Let $K\ge3$. If the color classes of the rows in $S$ and the
coalition $A$ are independent, a word omitting $S$ wins $r:A$
exactly when $A\cap T_r(S)\ne\varnothing$. For independent coalitions
$A_1,\ldots,A_d$ at $d\ge2$ distinct owners, a word omitting $S$
wins all these tests exactly when their intersection contains an
input outside $h(S)$, with no restriction on the row color classes.
\end{theorem}
\begin{proof}
For the second assertion such an input gives a projection. Conversely,
if every $v$ in the intersection is a head, retain one row at each
such $v$. For every $j$, the neighbors of $v$ miss $A_j$ since $v$
belongs to the independent set $A_j$. A word winning all the tests
cannot be atomic on any retained row. Indeed, for any proposed atom one of the required owners
is different, and its coalition misses the atom's row coalition,
contrary to common source (\cref{core:thm:profiles}). At the root both
children therefore omit the retained marks. One child keeps all the
demands: foreign wins pass to both and the gate owner's win, if
requested, passes to one. Descending with these same rows reaches a
projection in the intersection that fails its own retained row.

For the first assertion, if $A$ meets $T_r(S)$, fold that support at
$r$. For necessity, form $D_r$ in rounds $D^0=\varnothing$,
adding at round $d$ the heads of all rows with $t_r\subseteq D^{d-1}$.
For each new head choose one such row and let $S_d$ contain all chosen
rows through round $d$. We show that no word omitting $S_d$ wins any
independent coalition $A\subseteq D^d$ at $r$. At $d=0$ no word
wins the empty coalition. Otherwise take a shortest witness. A
projection in $A$ fails its own chosen row. If the word is atomic on
a chosen row at owner $r$, it wins that row's $r$-coalition, which
is independent and lies in $D^{d-1}$, while omitting $S_{d-1}$.
This contradicts induction. An atom at a different owner would give
two wins whose coalitions intersect only in $D^d=h(S_d)$, contrary
to the second assertion. With no atomic output on these rows both
children omit their marks and one still wins $r:A$, contradicting
shortestness. At stabilization this excludes every independent
$A\subseteq D_r$.
\end{proof}

Consequently, when all color classes are independent one can test
a mark $(i,\varnothing)$ in polynomial time, with no bound on color
repetitions. Start at $S=Z$. If head
$i$ is absent use its projection. Otherwise, whenever some $T_r(S)$
meets a remaining $r$-tail, delete all rows with such a meeting and
continue. The $r$-fold on this support omits $S$, absorbs on the
deleted rows and cannot absorb at the empty coloring. Attaching it
by $\star_r$ lifts any witness for the smaller set. If the process
stops with head $i$ still present, the theorem excludes every atom
on every remaining row for a word omitting $S$. Descent through a
child carrying the positive mark would then reach the forbidden
projection $i$. Each deletion is strict and each closure is finite
Horn propagation.

\begin{corollary}[Recognition with independent color classes]
\label{core:cor:triangle-free-marked-paths}
If every negative color class is independent, a query $(i,p)$ is decidable in
$2^{O(b\log(b+2))}$ times a polynomial in the explicit input size,
where $b=\max_r|p_r|$. More precisely, let $B_r$ consist of the
members of $p_r$ posted both at $r$ and at some other owner in the
negative rows. Put $L_r=\sum_{h=0}^{|B_r|}|B_r|!/(|B_r|-h)!$ and
$n=|V|$. At owner $r$
the path test of \cref{core:thm:palette-overlap} needs only lists of
distinct members of $B_r$ followed by a terminal input, hence at most
$nL_r$ paths. An impossible query has a witness of at most
$1+n\sum_r L_r$ negative rows.
\end{corollary}
\begin{proof}
For a gap at owner $r$ put $B=p_r$ and let $A$ be the demanded
$r$-tail. In the construction of
\cref{core:thm:pointed-guard-cut} an $r$-offer combines the two
smaller witnesses by $\star_r$. A foreign offer $s:C$ requires a
witness winning both $r:A$ and $s:C$ on $S_{s,C}$. By
\cref{core:thm:triangle-free-neighborhood-profile} some
$v\in A\cap C$ is absent from its heads. It is already absent from
$h(S)$, since a row headed by $v$ cannot post $v\in C$.
If $v\notin B$, its projection solves the entire gap. Otherwise
$v\in B_r$ by its occurrences in $A$ and $C$. Use
the unary operation $H_v(x)$ returning $\{s'\}$ when $v$ is any
foreign singleton $\{s'\}$, and $x+x_v$ otherwise. It is a core
term: successively apply $y\mapsto(y+x_v)\star_{s'}x_v$ for every
$s'\ne r$. Its $r$-profile on $D$ is $\Phi_r(x)(D)\land[v\in D]$.
Delete every row posting $v$ at any foreign owner. This includes all
rows covered by the offer. On the smaller set construct the original
gap witness and apply $H_v$: it absorbs all deleted rows, adds no
mark elsewhere since $v$ is head-free, and preserves both gap tests.
No descendant can select $v$ at another foreign offer. Induction gives
the claimed lists, with $r$-folds at their common prefixes.

Explicitly, the suffix support after $k$ letters consists of the
remaining letters and the terminal input. Discard paths whose full
support lies in $B$. On a row $(j,t)$ scan each path until its first
foreign singleton. An earlier letter $j$, or a reached terminal $j$,
emits the mark. For a first emission after $k$ letters on $w$ use the
Horn rule \eqref{core:eq:ordered-guard-horn}: discard $w$ if every
path $u$ whose suffix after $\min(k,[w,u])$ lies in $t_r$ is already
discarded, where $[w,u]$ is the common prefix length. The proof of
\cref{core:thm:palette-overlap} applies verbatim: each $H_v$ either
absorbs or adjoins one fixed menu, so shared prefixes merge and each
suffix support reads its $r$-profile. The word assembled from the
surviving paths wins on $A$ exactly when one has full support in $A$.
As there, repeatedly delete rows
with winning gaps until $i$ is head-free or no deletion is possible.
At a failing set keep a row headed by $i$ and one reason for each
noninitial path discard, reproducing the same failure and the stated
size bound. In the three-input example preceding
\cref{core:thm:repeated-blocker-budgets}, the color class $\{y,z\}$
contains an edge of the triangle. There the test already fails at
the empty coloring.
\end{proof}

\paragraph{A common theory of singleton tests.}
For arbitrary negative rows one can compute all positive queries from
a common theory. Choose a finite set of bits $d_{r,A}$, with owner $r$
and coalition $A\subseteq V$, including $d_{r,t_r}$ for every row
$(j,t)\in Z$ and $d_{r,p_r}$ for each query $(i,p)$ to be tested,
at every owner. Write $v_j(d_{r,A})=[j\in A]$ for the profile vector
of a projection. Let $\mathcal H$ consist of the clauses of
\cref{core:thm:profile-relations}: a positive unit, an implication
between two bits of the same owner, or a forbidden conjunction of bits
of distinct owners (including the empty conjunction). Write $H(P)$
for the clauses of $\mathcal H$ true at every vector in $P$, and
$C(\Sigma)$ for those entailed by a set of clauses $\Sigma$.

\begin{theorem}[A common Horn theory for source queries]
\label{core:thm:joined-profile-theory}
For $S\subseteq Z$, let $E(S)$ consist of the pairs $e=(r,A)$ with
$A\ne\varnothing$ occurring as $t_r$ in $S$. Write $d_e=d_{r,A}$
and $S_e=\{(j,t)\in S:A\not\subseteq t_r\}$. Put
\[
 B_S=C\biggl(\bigcup_{e\in E(S)}T_{S_e}\biggr),\qquad
 T_S=H(\{v_j:j\notin h(S)\})\ \cap\
 \bigcap_{e\in E(S)} C\bigl(B_S\cup\{d_e\}\bigr).
\]
When $E(S)$ is empty, only the projection theory remains; this includes
$S=\varnothing$. Every other $S_e$ in the formula is a proper subbudget.
The models of $T_S$ are exactly the profile vectors of words omitting
$S$. Without prescribing a profile, a word omitting $Z$ can supply
$(i,p)$ exactly when the following obstruction is absent:
\begin{equation}\label{core:eq:joined-profile-quiet}
 \exists S\subseteq Z:\qquad i\in h(S),\qquad
 (d_{r,t_r}\Rightarrow d_{r,p_r})\in T_S
       \quad((j,t)\in S,\ r\in I).
\end{equation}
More generally, prescribe an exact profile vector $x$ and any finite
set of queries to supply while omitting $Z$. This is possible exactly
when $x\models T_Z$, none of the requested queries is forced quiet,
and $x(d_{r,p_r})=0$ for each requested $(i,p)$ and every $r$.
\end{theorem}
\begin{proof}
The profile vectors attained on any budget are closed under all owner
gates: a gate on two words omitting a row also omits it.
By \cref{core:thm:profile-relations} these vectors are precisely the
models of their $\mathcal H$-theory. Induct on $|S|$. A model of $B_S$
is attainable separately on each lower budget $S_e$. If it makes
$d_e$ true for $e=(r,A)$, its witness on $S_e$ has value $\{r\}$
on every removed row, since $A\subseteq t_r$ there. Each of these
rows omits its mark, since the marked owner has no singleton input
there. Thus every model of a conditioned theory is attained on $S$.
Conversely every word omitting $S$ either has a singleton on one of its
rows, is an unmarked projection, or has two children that both omit
$S$: in the last case the root takes their union on every row.
The singleton case gives a true $d_{r,t_r}$ with $t_r\ne\varnothing$
and a model of the corresponding conditioned theory, as the word fits
every lower budget. When $E(S)$ is empty, no word is singleton on
any of its rows. Induction on the word
therefore shows that these models and the unmarked projections generate
all attained profile vectors. The theory of a generated union is the
intersection of the theories of its generators, giving $T_S$.

Suppose $S$ witnesses \eqref{core:eq:joined-profile-quiet}. A word
supplying $(i,p)$ is ambiguous there: its mark occurs only in the
head's nonsingleton menu. If the word omits $S$, the implications in
$T_S$ force it to be ambiguous on every row of $S$ as well. Both
children then omit $S$, and one supplies the requested mark. Descent
would reach the projection $i$, which a row of $S$ marks.
Conversely suppose no subbudget witnesses the display. On a current
budget $S$ with $i\notin h(S)$, projection $i$ works. Otherwise a
failed implication supplies an $S$-bounded word winning a test
$d_{r,t_r}$ and losing $d_{r,p_r}$. Hence $e=(r,t_r)\in E(S)$.
On $S_e$ a second word supplies the mark by induction.
Their $r$-choice absorbs on the deleted
rows, omits the others, and retains the requested mark.

For the last assertion, take a word $t$ of profile $x$ omitting $Z$.
If there are requested marks, take the union $w$ of their separate
bounded witnesses. As in \cref{core:prop:profile-maximum}, the reset
$w\del t$ is a core word: it returns $t$ on singleton outputs of $t$
and $w+t$ elsewhere. It keeps all profiles of $t$ and
omits $Z$. At each requested mark the zero profile bits make $t$
ambiguous, so the reset includes $w$ there. With no marks take $t$
itself. Conversely an attained mark precludes every singleton at its
query, and its word supplies the other two conditions directly.
\end{proof}

All entailments above are graph and subset tests. To test a clause
with positive premise $D$, start at its bits and all positive units
and follow the same-owner implication edges. If the reached set
contains a forbidden conjunction, the premise is inconsistent and
entails every conclusion. Otherwise the reached set itself is a
model, so it entails exactly its reached bits and does not imply
falsehood. With $b$ listed bits there are at most
$b+b^2+(b+1)^K$ clauses to test.
In particular, with no prescribed profiles or marks the omission
budget alone is feasible exactly when $T_Z$ lacks the empty conflict.
An inconsistent theory contains every implication, including those
in \eqref{core:eq:joined-profile-quiet}. A color absent at a query
has $p_r=\varnothing$, whose bit is false for every core word.

For all explicitly queried marks, the time at fixed $K$ is a polynomial
in the total input length times $\prod_r a_r$, where $a_r$ counts all
antichains, including the empty one, in the inclusion order on the
distinct nonempty $r$-coalitions of $Z$. Indeed, successive deletions
accumulate an upset at each owner. Its minimal elements form an
antichain and the upsets together determine the remaining budget.
Start the table at $Z$ and compute only its reachable subbudgets,
with polynomial work per budget for fixed $K$. All requested marks
use this table. For each, start at $S=Z$: if its head is absent the
projection succeeds; otherwise delete the full upset of any $e$
giving a failed implication in \eqref{core:eq:joined-profile-quiet}.
Restriction and the $r$-choice in the proof give equivalence before
and after this deletion. If all implications hold, $S$ certifies
failure. There are at most $m$ deletions per query. In the notation
of \cref{core:thm:pointed-guard-cut}, the antichain factor is at most
$(m+1)^W$, since an antichain at $r$ has at most $w_r$ elements.

The cut of \cref{core:thm:pointed-guard-cut} also has a quantified form involving only the
supplied coalitions, bits, and one input. Write $e=(r,A)$ for an offer
of a nonempty coalition occurring at $r$ in $Z$, and allow a blank
offer. An offer covers $(j,t)$ when $A\subseteq t_r$; a blank covers
nothing.

\begin{corollary}[An alternating transversal condition]
\label{core:cor:alternating-transversal}
Let $m=|Z|$. The positive $(i,p)$ is feasible exactly when
\[
  \exists e_1\ \forall b_1\in\{0,1\}\ \cdots\
  \exists e_m\ \forall b_m\in\{0,1\}\ \exists v\in V
\]
satisfies the following two conditions:
\begin{enumerate}
\item Every negative with head $v$ is covered by some $e_h$.
\item Retain those nonblank offers for which $b_h=1$, and let $a_r$
be the last retained coalition of each represented owner $r$. Then
$v\in L(a)$.
\end{enumerate}
An observer-agreeing antitone guard family with a positive bit is
realizable exactly when this condition holds at each maximal positive,
with its minimal zeros for $Z$.
\end{corollary}
\begin{proof}
After each offer delete the rows it covers. Bit zero keeps the previous
assignment, and bit one replaces its entry at the offered owner; a
blank does neither. The two successor tests in the proof of
\cref{core:thm:pointed-guard-cut} justify this step even if the
coalition occurs only in previously deleted rows: its own singleton
still enforces every row it covers. A successful construction can
choose an offer from a remaining row at each step until a projection
fits, and then use blanks. At most $m$ offers are needed.
Conversely the two conditions at the end say exactly that the
projection $v$ works on the remaining rows with the last retained
coalitions. Combining the two continuations at each nonblank offer
gives a witness. Union over maximal positives proves the final claim.
\end{proof}

The same game can read the prime guards directly. Allow any nonempty
coalition $A\subseteq V$ at any owner, with blanks as before. Use
at most $K(2^{|V|}-1)$ rounds. At the final input $v$ require
$v\in L(a)$ and, for every observer $k$ and passive coloring $t$
at $v$ omitting $k$,
\[
 \bigl(\text{no offered }(r,A)\text{ has }A\subseteq t_r\bigr)
 \quad\Longrightarrow\quad
 \text{some }G\in\Gamma_{k,v}\text{ is enabled at }t .
\]
Apply this test at every prescribed positive. The terminal implication
says exactly that every loss at $v$ was covered; covering all minimal
losses also covers their supersets. An arbitrary offer is sound by
the same two-child construction, since its second child passes the
offered singleton test. For completeness the original strategy
always offers an observed coalition deleting a remaining row,
so no pair $(r,A)$ recurs on a branch. This gives the round bound
without listing the minimal zeros. In particular different
presentations need not have comparable game lengths. If one head
is lost when its $r$-postings meet each of $d$ disjoint input pairs,
with all other heads retained, it has $d$ prime exclusions on that
observer facet and $2^d$ minimal losses. Against repeated declines,
offers restricted to observed tails must cover them one by one to
supply that head at the empty coloring. The two singleton offers
from either pair suffice in the unrestricted game: a saved offer
can terminate at its retained input.

In particular, recognition from the complete lists of minimal zeros is
in PSPACE at fixed $K$, even when the maximal positives are not supplied.
Enumerate one head and one proper passive coloring at a time. Inclusion
in one of the zero upsets decides its prescribed bit. Run the sample
test on every positive, reusing the workspace and requiring at least
one positive. This representation can prescribe exponentially many
positive bits from a short list.

\paragraph{A covering condition directly on the prime lists.}
The negative rows can be eliminated from the statement. Take a finite
rooted tree with edges labeled by pairs $(r,A)$, where $r$ is an owner
and $A$ a nonempty set of inputs. At a vertex $x$ write $E_x$ for
its outgoing labels. Siblings have distinct labels, and
$E_y\subseteq E_x$ at every child $y$ of $x$. Let $a_x$ give
the last coalition on the path to $x$ at each owner occurring there.
The available colorings at $x$ are those proper passive colorings
$t$ all of whose nonempty blocks occur as labels $(r,t_r)\in E_x$.
Availability at a head $v$ excludes postings of $v$, as usual.
At a leaf only the empty coloring is available.

\begin{theorem}[A prime guard cover on trees of coalitions]
\label{core:thm:prime-tree-cover}
Fix an observer-agreeing antitone family and a requested positive
$(i,p)$. A core word supplies this mark and respects all prescribed
zeros if and only if every tree as above has a vertex $x$ and an
input $v\in L(a_x)$ for which the prime guards at $v$ cover every
available coloring at $x$. Coverage means that a guard is enabled on
the coloring's observer facet. For a complete family with at least
one positive, testing this condition at its maximal positives
characterizes joint realizability.
\end{theorem}
\begin{proof}
If a tree violates the condition, choose at each vertex and each
eligible input one available coloring missed by its guards. For a
vertex $x$ let $S_x$ be all the resulting negative rows in its
subtree. Every eligible projection has a chosen loss in $S_x$. A shortest
feasible word must be singleton at a row of $S_x$, by the descent in
\cref{core:thm:pointed-guard-cut}. If its owner is $r$ and its block
is $A$, the label $(r,A)$ occurs where that row was chosen and hence
at $x$. The word itself meets the updated register at the
corresponding child, including the required gap when $a_x$ was
empty. It still omits all the child's rows. Induction from the
leaves rules this out.

Conversely start with an inclusion-minimal failing set of minimal
losses for the blank register. At a current minimal failing set
$S$ put all its nonempty blocks in $E_x$. Each eligible input
heads some row of $S$, whose coloring is available and missed.
For an outgoing $(r,A)$ delete the rows with $A\subseteq t_r$.
The old register is feasible on the proper subset left. The updated
register must fail there by the two-child construction. Use a
minimal failing subset for that child and repeat. This terminates
and gives the required tree. In fact a child reached through $(r,A)$
retains no label $(r,B)$ with $A\subseteq B$. Finally take the
union of the separate witnesses for the maximal positives.
\end{proof}

The trees in this condition can be large even for short prime lists.
For a cover crown, the same matching rows recur below many different
sequences of offers.

\begin{proposition}[Size of a tree witnessing noncoverage]
\label{core:prop:prime-tree-size}
For the opposite query of a cover crown of length $m$, every tree
witnessing failure in \cref{core:thm:prime-tree-cover} has at least
$2^m$ vertices, already counting only paths that offer the same owner
throughout. On $n$ inputs at fixed $K$, the worst-case minimum size
of a failing tree has double logarithm $n+o(n)$ to base two. The
lower rate holds for polynomial-size canonical prime lists.
\end{proposition}
\begin{proof}
Use the matching intervals $(L_j,U_j)$, indexed from $1$ to $m$
cyclically, and the opposite query $(L_1,U_m)$. Each matching
interval is a cover, so its only negative coloring has its head at
the sole free input and posts $c$ on $L_j$ and $b$ outside $U_j$.
These are all the negative colorings; more than one observer may
read the same row. The rows in any failing tree's root subtree
must include all $m$: deleting any one gives a patch supplying the
opposite mark (\cref{core:lem:crown-patch}). Its root therefore has
all $m$ outgoing labels $(c,L_j)$.

At a vertex reached using only $c$ offers and ending in $(c,L_j)$
the register asks to win on $L_j$ and lose on $L_1$. For every
$1\le k<j$, its subtree must contain matching row $k$. Otherwise
take the greatest bichain completion of the monotone Boolean function
\[
             f(X)=[\,L_h\subseteq X\text{ for some }k<h\le j\,].
\]
It wins on $L_j$ and loses on $L_1$, since the lower endpoints form
an antichain by \eqref{core:eq:crown-incidence}. It omits the mark at
every matching row except $k$: $f(U_\ell)$ reads membership of
$\ell$ or $\ell+1$ in the cyclic interval $(k,j]$, whereas
$f(L_\ell)$ reads membership of $\ell$ alone. This would realize the
register on all the subtree rows, contradicting the preceding proof.
Nested outgoing labels thus force a child $(c,L_k)$ for every $k<j$.
Every decreasing sequence of indices occurs as a distinct path from
the root, including the empty one. There are $2^m$.

For the upper bound put $d=K(2^n-1)$. In the tree constructed in the
theorem no edge label repeats along a path, so there are at most
$\sum_{h=0}^d d!/(d-h)!$ vertices. Its double logarithm is at most
$n+O_K(\log n)$. \Cref{core:cor:compact-locality-rate} gives the
matching lower rate using polynomial-size prime lists. For all
sufficiently large arities the cover crowns of
\cref{core:thm:interval-mask-locality-rate} give
$m\ge 2^{n-O(\log n)}$ without the prime-size restriction.
\end{proof}

Sharing subtrees changes this bound. Allow a finite rooted acyclic
directed graph, with the register written on each vertex and updated
along each edge. The same local test is exact, by unfolding the graph.
For a crown of length $m$ the least number of vertices in such a
certificate is exactly $(m+1)(m+2)/2$. Indeed, after a first offer $(c,L_j)$,
the preceding proof forces matching row $k$ somewhere below it for
every $k<j$. This forces edges $(b,V\setminus U_k)$ at that vertex
as well. Their children carry $\binom m2$ distinct two-owner
registers, in addition to the $2m$ forced one-owner registers and
the blank root.

For the upper bound take the path from $U_m$ to $L_1$ that reverses
matching rows $m,m-1,\ldots,1$ in order, using the inclusion
$L_j\subseteq U_{j-1}$ between reversals. Keep one vertex for each
contiguous segment of this path's list of reversals, including the
$m+1$ empty segments at its cuts. At a segment use the labels
$(c,L_j)$ and $(b,V\setminus U_j)$ of all its rows. The first edge
keeps the suffix after $j$ and the second keeps the prefix before $j$.
The endpoints $P,N$ of the segment start as $U_m,L_1$; a $c$-edge
replaces $P$ by $L_j$, and a $b$-edge replaces $N$ by $U_j$.
The segment boundaries thus determine both the register and its
eligible set $P\setminus N$, including the sole free head at the
root. Each $v\in P\setminus N$ must leave the running set somewhere
on this path. An inclusion cannot remove it, so it is the free head
of a reversal in the segment. That row supplies a missed coloring
using the vertex's outgoing labels. For an empty segment $P\subseteq N$.
This proves noncoverage everywhere with at most
$(m+1)(m+2)/2$ vertices.

\paragraph{Which lattices of remaining rows occur?}
Keep the rows left uncovered by a sequence of offers, irrespective of
the replies. These subsets of $Z$ form a finite lattice $L$ under
inclusion: concatenating two sequences takes the intersection, and
the empty sequence leaves $Z$. Write $z_r$ for the $r$-coalition of
row $z$. For $T\subseteq Z$ the least reachable set containing $T$ is
\begin{equation}\label{core:eq:offer-closure}
 \overline T=\{z\in Z:\ \forall r\text{ with }z_r\ne\varnothing,
                 \exists y\in T,\ z_r\subseteq y_r\}.
\end{equation}
Indeed, every offer that removes such a $z$ also removes a row of $T$.
For a row failing the test, offer its offending block $z_r$. Doing so
for every failing row leaves exactly $\overline T$. In particular,
allowing unobserved coalitions gives the same lattice. Each membership
$z\in\overline T$ uses at most $K-1$ witnesses in $T$.

For minimal losses this lattice distinguishes exactly the histories
that have different answers to a future head-coverage test. If $z$
survives in one history alone, let $j$ be its head. For every row
$y$ of the other history headed by $j$, some block $y_r$ is not
contained in $z_r$, by minimality. Offering those blocks covers
every remaining loss at $j$ in the other history and leaves $z$
in the first. Histories with identical surviving rows give the same
coverage after every continuation. The accepted coalitions in the
second terminal condition of \cref{core:cor:alternating-transversal}
are recorded separately.

Here is an intrinsic description of these lattices. In a finite
lattice write $J$ for the nonzero join-irreducibles and $M$ for the
nonunit meet-irreducibles. Put $M_j=\{m\in M:j\nleq m\}$.

\begin{theorem}[Owner colorings of the lattice of remaining rows]
\label{core:thm:offer-lattice}
Let $K\ge3$. A finite lattice occurs as the lattice of remaining
minimal losses of a core operation on $K$ owners if and only if there
is a coloring $c:M\to I$ such that, for every $j\in J$,
\[
 c(M_j)\ne I,\qquad
 j\nleq\bigvee\{m\in M_j:c(m)=r\}\quad(r\in I).
\]
Every such lattice occurs with one loss per conditional source and
all its postings on retained inputs.
\end{theorem}
\begin{proof}
In a lattice of remaining rows, each element $x$ is the join of the
closures of its individual rows. Thus for $j\in J$ choose $z$ with
$j=\overline{\{z\}}$: membership of $z$ in any $x$ means $j\le x$.
Every element is an intersection of single-offer survivor sets, so each
$m\in M$ is itself the survivor set of one offer. Color $m$ by the
owner of a chosen such offer. If $m\in M_j$, its offer removes $z$;
hence its color occurs on $z$, which omits an owner. For a nonempty
family of offers at one owner the union of their survivor sets is
itself a survivor set: offer the union of their coalitions. Thus this
union is their lattice join. If every one removes $z$, their join
still omits $z$. The empty join lies strictly below $j$. This proves
both conditions.

Conversely represent $x$ by $J(x)=\{j\in J:j\le x\}$. Make one
head for each $j$ and one retained input $u_m$ for each $m\in M$.
At head $j$ prescribe the single minimal loss posting exactly
$u_m:c(m)$ for $m\in M_j$. This coloring is proper by the first
condition; it is nonempty since the meet-irreducibles separate $j$
from the bottom element. These losses occur together in a core
operation: for each head take its one-source completion on the common
retained inputs (\cref{core:thm:source-completion}) and take their
union. A head absent from a summand has no source there. For the
one-element lattice use a projection with no losses.

An offer can remove a row only if it consists of retained inputs
$u_m$ all posted at its owner, say with nonempty index set $X$. Its
survivors are $\bigcup_{m\in X}J(m)$. This equals $J(\bigvee X)$:
a join-irreducible absent from every $J(m)$ is absent from their join
by the second condition. Ineffective offers leave the top. Every
history therefore leaves a set $J(x)$, and offering $u_m$ alone at
$c(m)$ leaves $J(m)$. Since every $x$ is a meet of meet-irreducibles,
all $J(x)$ occur. The order on them is the original lattice order.
\end{proof}

For example, take a height-two lattice with $d\ge3$ atoms. Two
different atoms join above any third, so the atoms need distinct
colors. Conversely $d$ colors suffice, since $M_j$ omits $j$.
Adjoining a new top above this lattice raises the least number of
owners to $d+1$: the new top is join-irreducible and sees all $d$
colors among the atoms it misses. Giving the old top any one of
those colors and leaving the extra color unused meets both conditions.
Thus even among acyclic realizable families the number of owners
gives a strict hierarchy of lattices. A lattice type by itself cannot
exclude a candidate family: every type arising from proper rows
already occurs among realizable families, by the same construction.

For three and four owners the coloring test is ordinary graph
colorability. Put vertices $J\sqcup M$, join $j$ to each member of
$M_j$, and join $m,m'$ if their join lies above a $j$ missing from
both. A lattice coloring extends to this graph by assigning each
$j$ a missing color. Conversely suppose the graph is properly colored.
The maximal members of $M_j$ have distinct colors: the join of two of
them lies above $j$ by maximality. Consider a color class $X$ in
$M_j$. If its color occurs at a maximal member $m$, every member of
$X$ lies below $m$, by the same edge test. Otherwise at most $K-2$
maximal members occur: both the color of $j$ and the color of $X$
are missing among them. If no maximal member bounds $X$, one can
choose a member of $X$ outside each such bound. Their join lies above
$j$, since any element missing $j$ lies below a maximal member of
$M_j$. For $K\le4$ this contradicts the pair test. The same argument
for general $K$ shows that monochromatic subsets of size at most
$\max(2,K-2)$ suffice in \cref{core:thm:offer-lattice}.

\paragraph{An implicational test and the size of its countermodels.}
There is also a test in ordinary intuitionistic implication. It uses
one atom $d_z$ per listed negative and one atom $w_a$ per register
assignment, including the blank assignment $\varnothing$. For a list
$B$ of rows let $[B]w$ abbreviate the nested implications with
premises $d_z$ for $z\in B$ and conclusion $w$; an empty list gives
$w$. Write $Z_v$ for the rows with head $v$ and $C_e$ for those
covered by a nonblank observed offer $e$. Use the hypotheses
\[
 [Z_v]w_a \quad(v\in L(a)),\qquad
 ([C_e]w_a)\ \to\ ([C_e]w_{a[e]})\ \to\ w_a
       \quad\text{(every $a,e$)},
\]
where $a[e]$ replaces the entry at the offered owner.
These formulas have polynomial total length on an explicit sample at
fixed $K$: there are at most $(1+K|Z|)^K$ assignments.

\begin{proposition}[An implicational criterion for guard realization]
\label{core:prop:guard-sequent}
The hypotheses entail $w_\varnothing$ in intuitionistic logic exactly
when the marked query is feasible. On failure the lattice of remaining
rows, ordered by reverse inclusion, supplies a finite Kripke
countermodel. If every failing subfamily of the negatives has at
least $b$ rows, then every finite Kripke countermodel has width at
least $\lceil b/|V|\rceil$.
\end{proposition}
\begin{proof}
Read a winning tree from its leaves upward. At a leaf use the first
hypothesis for its eligible input; each required $d_z$ has already
been supplied on the path. At an offer abstract the assumptions for
$C_e$ in both continuation proofs and apply the second hypothesis.
This gives a derivation of the root atom with no row assumptions.

On the reverse lattice force $d_z$ at $S$ when $z\notin S$, and
$w_a$ when $a$ is feasible on $S$. Both predicates persist upon
deletion. The first world above $S$ forcing every $d_z$ for $z\in C_e$
is $S\setminus C_e$. Thus $[C_e]w_a$ holds at $S$ exactly when
$a$ is feasible on this smaller set. The two-child construction of
\cref{core:thm:pointed-guard-cut} validates the offer hypothesis;
a fitting projection validates the leaf hypothesis. Failure makes
$w_\varnothing$ false at the initial world $Z$.

In any finite countermodel of width $q$ consider only worlds above
and including the refuting root. For each head $v$ the worlds where
at least one $d_z$ for $z\in Z_v$ is false form a downset.
At each of its maximal worlds retain one such row. At most $q|V|$
rows are retained. Persistence means that a false premise still
exists everywhere in that downset. The leaf hypotheses therefore
still hold: whenever $w_a$ is false, every eligible head has a false
premise at that world. Removing premises from $[C_e]w$ makes this
antecedent of the offer rule harder to satisfy, so all remaining offer
hypotheses still hold. Discard the offers and register assignments
absent from the restricted sample. The same root still refutes the
conclusion, so the retained sample fails, proving the bound.
\end{proof}

The implication test has a succinct form on the canonical primes
alone. Use finite, constant index domains for the inputs, owners,
proper passive colorings, nonblank offers and register assignments;
their indices have polynomial bit length. Let $X_r(A)$ be an atomic
predicate on nonempty coalitions, subject to
$X_r(A)\to X_r(B)$ for $A\subseteq B$.
For a coloring $t$ at $v$ write $P_v(t)$ for the rigid input test
that some listed prime guard is enabled on its observer facet.
Agreement makes the choice of a missing observer immaterial. Put
\[
 Q(v)=\forall t\text{ at }v\quad
 \left(P_v(t)\ \lor\!
       \bigvee_{r:\,t_r\ne\varnothing} X_r(t_r)\right).
\]
Use just the following two universally quantified hypotheses, in
addition to monotonicity:
\[
  Q(v)\to W(a)\quad(v\in L(a)),\qquad
  (X_r(A)\to W(a))\to(X_r(A)\to W(a[r:=A]))\to W(a).
\]
The query succeeds exactly when these hypotheses entail
$W(\varnothing)$. Quantifiers here abbreviate finite conjunctions or
finite case distinctions over the stated domains. Keeping the indices
bound gives polynomial length in the prime presentation; expanding
them can take exponential space.

Indeed, along a successful branch the assumed offers cover every
losing coloring of the terminal input. For each such coloring,
monotonicity gives $X_r(t_r)$ from a covering offer, proving $Q(v)$
by cases. Conversely use the same countermodel with $X_r(A)$ true
at $S$ when no row left in $S$ is covered by $(r,A)$. This predicate
is monotone in $A$ and persistent under deletion. The first extension
making it true deletes exactly those rows. If a loss at $v$ survives,
its own coloring falsifies $Q(v)$, since every nonempty block covers
that very row. If none survives, every losing coloring extends a
minimal loss at $v$ already deleted by some offer. That offer and
all enlargements of its coalition are true at $S$. Thus $Q(v)$
holds exactly when the leaf $v$ is available. The two hypotheses
follow as in the explicit sequent. Every countermodel of the indexed
form also gives one for the explicit sequent on the same frame: set
$d_z=\bigvee_{r:z_r\ne\varnothing}X_r(z_r)$, with the empty
disjunction false. Truth of all minimal-loss premises for $v$ implies
$Q(v)$ by monotonicity. Forcing an offer forces $d_z$ for every row
it covers, giving the explicit offer hypothesis.

\begin{proposition}[Shared certificates and finite countermodels]
\label{core:prop:guard-dag-countermodel}
For a failing indexed query, let $D$ be the least number
of vertices in a graph certifying noncoverage and $F$ the least number
of worlds in a countermodel. For $n=|V|$ one has
\[
                         F\le D\le F(1+nF)^K .
\]
\end{proposition}
\begin{proof}
For the first inequality order the graph by descent from its root.
Choose one missed row for each eligible input at each vertex.
Force $X_r(A)$ at $x$ when $(r,A)$ covers no chosen row at or below
$x$. Make $W(a)$ false at $x$ when some descendant $y$, possibly
$x$, has a register on the same owners containing $a$ coordinatewise.
The eligible inputs of $a$ are then eligible at $y$, whose chosen
rows refute their $Q(v)$ at $x$. To check the offer rule for
$e=(r,A)$, start at this $y$. If $e$ covers no chosen row below it,
the old register is still false there and $X_r(A)$ holds.
Otherwise some such row has a block $B\supseteq A$ at $r$.
Heredity gives an edge $(r,B)$ here. Take this edge and repeat until
$e$ covers no remaining row. The graph is finite and acyclic, so the
process terminates. The final register contains $a[r:=A]$, since
only $r$ was updated. Thus at least one of the two antecedents in
the offer rule fails at $x$. Both predicates persist under descent,
and the blank register fails at the root.

Conversely, use pairs $(x,a)$ where $x$ is a maximal world at
which $W(a)$ is false. Start at such a world for the blank register
above the refuting root. For each input $v$
choose a missed coloring at each maximal world failing $Q(v)$.
Its nonempty blocks have false atoms there, since a coloring
witnessing failure at an extension also witnesses it at the original
world. At $(x,a)$ use as labels the blocks $(r,A)$ of the selected colorings
for which $X_r(A)$ is false at $x$.
Every eligible $v$ at $x$ still has a selected coloring all of whose
blocks are false at $x$, by persistence. For each outgoing label $e$, failure of an
offer antecedent
gives a world above $x$ forcing its atom where the old or updated
register fails. This world is strictly above $x$, so maximality
excludes the old register. Extend to a maximal world failing the
updated register and use that pair as the child. Persistence gives
nested labels; monotonicity removes the whole upset of the accepted
coalition. At frame width $u$ at most $nu$ colorings were selected;
there are at most $1+nu$ register entries per owner, including absence.
This gives at most $F(1+nu)^K$ pairs and hence the bound.
\end{proof}

For the opposite query of a cover crown of length $m$ from
\cref{core:thm:exponential-locality}, the least possible countermodel
height is exactly $m+1$, counting worlds. As in \cref{core:prop:prime-tree-size}, write $z_j$ for
matching row $j$, use $(L_1,U_m)$ as the opposite query, and offer
$(c,L_j)$ in order $m,m-1,\ldots,1$. Each offer covers just $z_j$,
by \eqref{core:eq:crown-incidence}. At the root $d_{z_m}$ would
force $w_\varnothing$: deleting $z_m$ gives a patch supplying the
mark. The offer hypothesis for the first step therefore forces a strict successor
where $d_{z_m}$ holds and the register $c:L_m$ fails.

Consider the subsequent step from $c:L_j$ to $c:L_{j-1}$ for $j>1$.
Once $d_{z_{j-1}}$ also holds, the old register must succeed.
Indeed its task is to win on $L_j$ and lose on $L_1$ at $c$.
The Boolean function $X\mapsto[L_j\subseteq X]$ does this, and its
greatest bichain completion omits every row except possibly $z_j$
and $z_{j-1}$: these are the only upper endpoints containing $L_j$.
The derivation above, starting at this register, thus forces its atom
under the two row assumptions. Consequently $d_{z_{j-1}}$ is still
false at the present world. By the offer hypothesis it becomes true
at a strict successor where the new register fails. This gives $m$
strict steps. The reverse lattice on $m$ rows has height at most
$m+1$, proving equality. Every deletion patch supplies the opposite
query, so the width bound is at least $\lceil m/|V|\rceil$ as well.

The same crowns require $\Theta(m^2)$ worlds in total. Offer
$(c,L_j)$ and then $(b,V\setminus U_k)$ for $1\le k<j\le m$. Each forces a successor as above:
the deletion patch at $j$ forces the blank register under $d_{z_j}$,
and the separator for $(k,j]$ in \cref{core:prop:prime-tree-size}
forces the $c:L_j$ register under $d_{z_k}$. At the resulting world
$d_{z_k},d_{z_j}$ are true and the two-owner register fails.
Every $d_{z_h}$ with $k<h<j$ must be false there: the separator
for $(h,j]$ wins on $L_j$, is zero on $U_k$, and omits every mark
except $h$. Keeping $k\le\lfloor m/2\rfloor<j$, the nearest true
indices on either side of this cut therefore recover $k,j$.
This gives at least $\lfloor m^2/4\rfloor$ distinct worlds.
Pairs with both endpoints strictly increasing give
$\lfloor m/2\rfloor$ incomparable worlds, since each forces a row
false at the other. The interval graph gives a countermodel of
quadratic size by the preceding construction and has width at most
$m+1$: segments with a common left boundary are linearly ordered by
reachability (a prefix edge cuts at any row).
Thus minimum width has order $m$, and minimum height is $m+1$.
The lower bounds hold for the explicit as well as the indexed
sequent, by the same row implications. For the polynomial prime
lists of \cref{core:cor:compact-locality-rate} both dimensions can
be $2^{|V|-o(|V|)}$. At width four of the primes both already grow
exponentially (\cref{core:thm:compact-prime-locality}).


\paragraph{Offers fixed before any responses.}
A fixed list has a direct algebraic interpretation. A greatest profile
completion on a nonempty support $H\subseteq V$ has the prescribed
singleton outputs there and returns the union of the menus in $H$
whenever its output is ambiguous (\cref{core:prop:profile-maximum}).
Inputs outside $H$ are ignored.

\begin{theorem}[Fixed lists and greatest profile completions]
\label{core:thm:fixed-offer-completion}
A fixed list wins at $(i,p)$ against $Z$ if and only if a greatest
profile completion on some $H\ni i$ supplies that mark and omits $Z$.
Consequently the complete nonzero families admitting a fixed winning
list at every maximal positive are exactly the finite nonempty unions
of these completions, allowing different supports.
\end{theorem}
\begin{proof}
For a winning list let $H$ consist of the heads at which it covers
every negative, including inputs with no negatives. Declining every
offer puts $i$ in $H$. Accepting only $(r,A)$ gives
$(A\cap H)\setminus p_r\ne\varnothing$. Accepting one offer at each
of several distinct owners shows that their coalitions meet in $H$.
These conditions also suffice for every response, since only the last
acceptance at each owner counts.

On $H$ give $r$ the upward closure of its offered coalitions $A\cap H$,
adjoining $H$ at every owner. Every rainbow selection intersects by
the preceding tests, so \cref{core:thm:profiles} realizes the
profiles. None wins on $p_r\cap H$, since even $H$ contains the
unposted input $i$. The greatest completion supplies the mark.
For a negative headed in $H$, some covering offer makes its owner
the singleton output. A negative headed outside $H$ has no source.

Conversely, offer the minimal winning coalitions of the profiles on
$H$. A loss headed there must be atomic, since the ambiguous value
would contain the head's mark. Its winning coalition gives a covering
offer. All singleton tests lose at the positive, and the rainbow
condition gives every required intersection in $H$. This list wins.
With explicitly given negatives one can keep at most one covering
offer per row headed in $H$, enlarging it to that row's observed
coalition. All the required intersections persist. The complete-family
statement follows by taking the union of the positive witnesses;
conversely each summand respects all zeros and some summand supplies
each positive.
\end{proof}

Three inputs and two losses can already require adaptation. Use
the word $(y\star_c x)\star_b z$. Require $x$ quiet with $y,z$
posted at $b$, and $y$ quiet with $x$ posted at $c$. Ask for $y$
with $x$ posted at $b$ and $z$ at $c$. The outer and inner gates
respect the two losses, respectively, and both pass the mark at the
positive. To cover the loss at $y$ a fixed list must offer
$(c,\{x\})$. Accepting only this forces the loss at $x$ to be
covered as well, by a nonempty $b$-coalition in $\{y,z\}$. Accepting
both leaves no eligible input. Taking the word's complete guards
therefore gives a core operation outside the union class of
\cref{core:thm:fixed-offer-completion}. Two inputs cannot give this
separation: by \cref{core:prop:binary} every binary core operation
is its greatest completion on its support. One loss cannot give it
either: use projection $i$ if its head differs from $i$; otherwise
feasibility requires some posting $j\in t_r\setminus p_r$, and the
single offer $(r,\{j\})$ works.

When all negatives have the queried head $i$, nesting has an intrinsic
description that allows repetitions of any color. For a positive $(i,p)$
label each negative $(i,t)$ with a color $r$ satisfying
$t_r\setminus p_r\ne\varnothing$. A block of labeled rows is admissible
if every choice of one tail from each of several distinct labels has
nonempty intersection.

\begin{theorem}[Completion depth at one source]\label{core:thm:one-source-depth}
The minimum depth is the least number of admissible blocks partitioning
the negatives, minimized over their labels. The empty family has depth
zero; a row with no eligible label makes the query impossible. For
$n$ inputs, the maximum finite depth is $\min(K,n-1)$, including for
complete families with every source other than $i$ always retained.
\end{theorem}
\begin{proof}
Follow a child supplying the positive mark, starting at the root of a
fitting expression. A free union passes all remaining omissions to
that child. At a completion gate, assign each remaining negative on
which the output is a singleton to the block of this gate, labeled
by that owner $r$. The profile of the entire subexpression wins on
$t_r$ and loses on $p_r$, since it supplies the mark at the positive.
Its rainbow condition gives every required intersection. On all other
remaining negatives the gate is ambiguous and therefore passes their
omissions to the chosen child. At a leaf supplying the mark only
$x_i$ is possible, and it omits no negative headed at $i$.

Conversely, start with $x_i$ and attach the blocks in any order.
For one block take a completion on the original input ports. At each
owner generate its profile by the tails labeled with that owner and
by the full port set. Admissibility gives the rainbow condition. Plug
the previous expression into port $i$, leaving the other ports as
projections. Each labeled tail still forces its singleton, since it
omits $i$. On a previously enforced negative all ports omit the mark,
so the new gate preserves the omission. At the positive every labeled tail escapes $p_r$, and the full
port set escapes it because port $i$ is ambiguous. The gate therefore
supplies the mark from that port. Thus one layer per block suffices.

Grouping by eligible label gives at most $K$ blocks. Alternatively
choose a witness $v\in t_r\setminus p_r$ per row and group by $v$,
giving at most $n-1$. For sharpness take $\min(K,n-1)$ distinct
passive inputs with distinct colors, each a sole posting on its own
row, and the empty positive. No two rows fit in one block.
The complete family retaining every passive input and suppressing $i$
exactly above those rows has the same depth: take the free union of
the bounded witnesses for its maximal positives and all passive
projections. The same construction gives the upper bound for every
complete family retaining all other inputs.
\end{proof}

For single-posting losses with a common head, the exact number is the
minimum vertex cover size in the bipartite graph of their input--color
pairs, provided all postings are eligible. Within a block, either all
edges use one color, or the intersection test forces every edge to use
the very same input. Conversely every such star is admissible. With
larger tails even the depth-one test depends on the number of colors.

\begin{proposition}[One source at depth one]\label{core:prop:common-head-depth-one}
Suppose every negative has head $i$ and the positive is $(i,\varnothing)$.
With at most two posting colors, feasibility by a greatest profile
completion is decidable by 2-SAT. With three posting colors it is
NP-complete, even when each row has exactly two colors and a
depth-two realization always exists.
\end{proposition}
\begin{proof}
By \cref{core:thm:one-source-depth} one must choose a nonempty
tail per row so that every rainbow selection intersects. With two colors give each row a Boolean choice, fix
the one-option rows, and forbid each disjoint pair of different colors.
These are all the constraints. With three colors a guessed selection
is checked on at most triples, giving NP membership.

Reduce 3-SAT. Give each literal occurrence its own two options for
the underlying variable's truth value, in distinct colors. Within
each clause color its three falsifying options with three distinct
colors and forbid that triple. Assign any different color to each
remaining option. To enforce equal truth values at occurrences of a
variable, connect their ordered pairs of true and false colors by
paths of two-option rows, changing one color at a time. Such paths
exist: the six ordered pairs are arranged in the cycle
$12,13,23,21,31,32,12$. Between adjacent rows forbid the two
mismatched truth-value pairs. Each is rainbow, since the ordered
color pairs agree in one coordinate and neither repeats a color.
These constraints admit a selection exactly for satisfiable formulas.

Realize the forbidden pairs and triples as empty intersections. For
each rainbow subset of at most three options with at most one from
each row and containing no forbidden set, introduce one passive input.
Post it exactly at the options in that subset. The two tails of a
row are disjoint and nonempty. Any rainbow selection with at most
one option per row has a common input exactly when it contains no
forbidden set. This uses polynomially
many inputs. Add the common head $i$ without posting it; each
two-color row is proper. Finally fix two of the three colors.
Starting with $x_i$, fold the current word with all passive inputs
at each of those owners in succession. Each row posts at least one
chosen color, which absorbs there, and later passive inputs cannot
restore its head's mark. Both folds pass the mark at the empty
coloring, proving the depth-two bound.
\end{proof}

\begin{lemma}[Increasing the least sample depth by one]\label{core:lem:depth-increment}
Take a feasible sample with one positive and add a fresh input $y$,
unposted on the old queries. At every old input $v$ add two negatives
headed at $v$, with sole posting $y:b$ and $y:c$, respectively, for
two distinct colors $b,c$. The new minimum depth is the old one plus one.
\end{lemma}
\begin{proof}
The binary gate of \cref{core:thm:completion-suspension} fits the new
sample when applied to an old witness and $y$. Conversely take any
outer summand supplying the old positive. For $s=b,c$ its $s$-profile
must win on $\{y\}$. Indeed give $y$ its $s$ singleton and all old
inputs their $r$ singleton for some $r\ne s$. At observer $r$ the only
visible posting is $y:s$, so the new negatives exclude $r$ from
every old source. Only $s$ remains. By singleton-profile absorption
the summand contains $y$ everywhere. Thus on old probes it is
ambiguous. It cannot be a projection: only $y$ satisfies these two
profile tests, and $y$ omits the positive mark. A child of the outer
gate supplies that mark and omits every old negative. Erasing $y$ at
neutral status preserves those source bits by
\cref{core:thm:completion-erasure}, giving an old witness one layer lower.
\end{proof}

For every fixed $d\ge1$ and $K\ge3$, deciding depth at most $d$ is
NP-complete even for one positive at the empty coloring, at most two
posting colors per negative, and the promise that the optimum is $d$
or $d+1$. Start with the depth-one versus depth-two instances in
\cref{core:prop:common-head-depth-one}. Apply
\cref{core:lem:depth-increment} $d-1$ times. Each iteration adds one input and at most
twice the current input count in rows. Membership follows from
\cref{core:thm:layer-interpolation}.

\begin{corollary}[Two-color recognition of unions of profile completions]
\label{core:cor:fixed-offer-prime-space}
Fix $K\ge3$. When all minimal losses use only $b,c$, recognition of
unions of greatest profile completions from the complete prime guards
is PSPACE-complete, even on the directed-cube masks of
\cref{core:prop:compact-cube-masks}.
\end{corollary}
\begin{proof}
For a positive $(i,p)$ choose $H\ni i$ and put
$L=p_c\cap H$, $U=H\setminus p_b$. On $2^H$ take inclusion and the
reversals
\[
     H\setminus t_b\longrightarrow H\cap t_c
       \qquad ((j,t)\text{ a loss using }b,c,\ j\in H).
\]
There is a fixed winning list exactly when some $H$ has no path
from $U$ to $L$. Indeed, in the fixed-list proof other owners cover
no minimal losses and may be dropped. On its covered heads let $f$
be the monotone Boolean function with terms $A\cap H$ for the
$c$-offers and one more term $U$. Then $f(U)=1$, $f(L)=0$.
A $c$-offer covering a row makes $f(H\cap t_c)=1$.
A covering $b$-offer meets every term, including $U$, so
$f(H\setminus t_b)=0$. Thus $f$ respects every reversal.
Conversely, without the path, the vertices reachable from $U$
form an upset separating $U$ from $L$ and constant on each reversed
interval. Offer its minimal true sets at $c$ and the complements in
$H$ of its maximal false sets at $b$. Constancy covers every loss
headed in $H$, and separation and monotonicity give all the required
intersections.

Require at least one positive, then stream the positives and the subsets $H$. Testing an edge needs
only a scan of two-color partial colorings and prime evaluations.
Midpoint reachability on at most $2^{|H|}$ vertices uses polynomial
space. For hardness use the directed-cube construction of
\cref{core:prop:compact-cube-masks}. A cyclic instance is outside
the core. An acyclic mask is path-closed, and the proof of
\cref{core:thm:interval-masks} already expresses every such mask as
a union of greatest completions on $V$.
\end{proof}

\paragraph{The cost of one layer from prime lists.}
A prime list can prescribe exponentially many losses even when every
source reads only one earlier group of inputs. With three posting
colors, choosing compatible singleton outputs on these losses encodes
the truth tables of Boolean functions.

\begin{theorem}[Succinct recognition at prescribed depth]
\label{core:thm:succinct-layer-hardness}
For every fixed $K\ge3$ and $d\ge1$, deciding from complete canonical
prime guards whether depth at most $d$ suffices is NEXP-complete.
Hardness holds on realizable families of minimum depth $d$ or $d+1$
with acyclic source dependence. For $d=1$ every dependency path has
at most two edges.
\end{theorem}
\begin{proof}
Expanding all probes and applying \cref{core:thm:layer-interpolation}
gives nondeterministic exponential time. We prove hardness at depth one.

\emph{Addressed choices.}
Start with a dependency quantified Boolean formula. Each existential
variable chooses a Boolean function of its specified universal bits,
and the matrix must hold on every universal assignment. This problem
is NEXP-hard \cite{ChenHuangLuTan2022}. Use clauses of width three by
introducing gate variables depending on all universal bits. Remove
tautologies. A remaining all-universal clause makes the instance false
and can be replaced, together with the rest of the instance, by the
fixed contradiction $z\land\neg z$. A fresh satisfiable unit on its
own variable ensures there is at least one type.

A type has one row per bit string on its address coordinates. Each row
must choose one of two truth values with distinct colors among $a,b,c$,
or a single forced option. Make one master type per existential
variable, addressed by its dependency bits. Make a separate type for
each existential occurrence, with the same address coordinates. Repeat
an occurrence with fresh copies to give each clause three existential
slots when it has fewer. Give the three falsifying options distinct
colors. Forbid the triple of these options precisely when the addresses
agree on shared coordinates and agree with the values falsifying the
universal literals in that clause. These conditions characterize
extension to a common falsifying universal assignment.

Tie every occurrence to its master as in
\cref{core:prop:common-head-depth-one}. Join their ordered pairs of
true and false colors by a path changing one coordinate at a time;
a direct link also works for identical pairs. Use fresh intermediate
types on the same address coordinates. For each link forbid both
truth-value mismatches at equal addresses. Each mismatch uses two
colors. Adjoin a fresh one-option type of the third color with empty
address to turn it into a forbidden triple. Thus all forbidden patterns
have three distinct types and colors. Each triple of option types has
at most one prohibition, conditioned solely on equalities between
bits of different addresses and on fixed bits. A selection avoids all
prohibitions exactly when the input formula is true.

\emph{Intersections on polynomially many inputs.}
Construct a set $W$ so that a rainbow selection of at most three options
from distinct addressed rows has intersecting tails exactly when it is
allowed. For each nonempty rainbow pattern of option types with
distinct underlying types, allocate a private coordinate if it has no
prohibition, posting it at all its named options. If it has a
prohibition, allocate one coordinate for each way of violating one
address condition: the wrong value of a fixed bit, or one of the two
mismatches of equal bits. Post this coordinate at each named option
whose part of that witness holds, unconditionally at a named option
with no part to check. It belongs to all three tails precisely on that
mismatch. No coordinate is posted at an unnamed option. For a pattern
with two options of the same type, use coordinates witnessing a
differing address bit in each of the two orientations. Use the indicated
bit at the two options and post unconditionally at the third, if present.
The opposite options of a single addressed row consequently have
disjoint tails. A coordinate for a smaller pattern cannot fill a
forbidden triple.

For every type, address bit, and ordered pair of distinct colors of
its options, also add a private coordinate switching between the first
color at bit zero and the second at bit one. These serve only one type
and preserve all forbidden triples. In particular each option has an
always-posted private coordinate from its one-option pattern. There
are polynomially many coordinates. At each type any coordinate is
fixed or its posting depends on just one address bit.

\emph{The complete prime lists.}
Use $W$ as passive inputs and a disjoint copy $P$ of $W$. Copy every
$b$- and $c$-posting onto its twin in $P$. Give every $a$-option
one more passive input $m$, posted at $a$. Each type $\lambda$ has a
head $j_\lambda$ whose losses are the upsets of its addressed rows
with these enlarged tails. Add a head $i$. Its losses each post all of
$P$ at $a$ together with one further input at $a$ from
$J=\{m\}\cup\{j_\lambda\}$. All passive inputs are always retained.

Here these specifications give the full canonical prime lists in
polynomial time. If an observer is an option color of a type, every
addressed row uses that color and the facet has the empty guard.
On any other facet the posting set at address $x$ has the form
$F\cup\bigcup_h B_{h,x_h}$, with $F$ fixed and one block per address
bit. A guard hits every such row precisely when it hits $F$ or hits
both blocks of a single bit: otherwise choose an unhit block at
each bit. Every addressed row itself is a proper consistent coloring.
The minimal hitting sets are obtained by listing the singletons from
$F$ and the pairs from opposite blocks, then minimizing. At $i$ the
$a$-facet has the empty guard. On another facet exclude one $a$-posting
from $P$, or exclude all the $a$-postings from $J$. These are all the
primes. Dependence runs from passive inputs to type heads and from
passive inputs and type heads to $i$.

\emph{The distinguished positive.}
Ask for $i$ with just $P$ posted at $a$. In a depth-one realization
some completion on a support $H$ supplies this mark and fits all losses,
by \cref{core:thm:fixed-offer-completion}. It uses $i$. At an $i$-loss
its ambiguous value would pass the mark, so its $a$-profile wins on
$H\cap(P\cup\{j\})$ for every $j\in J$ and loses on $H\cap P$ at
the positive. Hence $J\subseteq H$. Every addressed row of every type
must now yield a singleton, as its head is in $H$. Choose that option
in the row. All rainbow selections intersect by admissibility. A
forbidden triple could not acquire an intersection on $P$ or at $m$:
no $a$-option posts in $P$, and only $a$-options post $m$. The original
choices therefore satisfy the formula.

Conversely take a satisfying selection. On full support generate
profiles by its chosen tails, the $a$-sets $P\cup\{j\}$ for $j\in J$,
and unanimity. Chosen options intersect on $W$ for every rainbow
selection. When a new $a$-set is selected instead, the other chosen
tails have a common coordinate in $P$ by the copying rule. Thus these
profiles are admissible, and their greatest completion omits every
specified loss. It supplies $i$ on every positive coloring posting all
$P$ at $a$: no member of $J$ then posts at $a$, every chosen
$a$-tail contains $m$, and every other chosen tail meets $P$.
If some $v\in P$ is missing at $a$, the term $x_i\star_a x_v$
already supplies this positive and fits all losses.

\emph{The other positives and the depth promise.}
At a positive headed at $j_\lambda$, either a fixed posting of that
type is absent or each block of some bit has an absent posting. In
the first case attach its input at its owner to $x_{j_\lambda}$.
In the second choose one absent posting from each block. If their
colors agree the two corresponding singleton coalitions are admissible.
If the colors differ, adjoin their private switching coordinate to
both coalitions. They now intersect; each still escapes the positive
and wins on its whole branch of addresses. In both cases the resulting
completion uses just $j_\lambda$ and passive inputs, omits all the
losses and supplies the requested mark. Passive heads use projections.
A free union of these witnesses completes the depth-one realization.

Regardless of satisfiability, two layers suffice. Only the $i$-positives
with all $P$ posted at $a$ need another witness. For a type with an
option $s\ne a$, choose its always-posted private input $v$ in $P$ and
use the helper $x_{j_\lambda}\star_s x_v$. For an $a$-only type use
$x_{j_\lambda}\star_a x_m$. Each helper omits every type loss, including
those at its own head. Take the $a$-fold of $x_i$, $x_m$ and all helpers.
At an $i$-loss its extra posting forces an $a$ singleton, through the
associated helper or through $x_m$. At the positive no helper is an
$a$ singleton since no member of $J$ posts at $a$. The fold passes
$i$ and introduces no lost mark. The optimum cannot be zero since $i$
is conditional. This proves the depth-one promise. Finally apply
\cref{core:thm:completion-suspension} $d-1$ times. Its prime-list
transformation is polynomial and each fresh input precedes the old
heads, preserving acyclicity.
\end{proof}


\paragraph{Depth of compositions.}
Count one layer for a greatest profile completion on any nonempty
support, applied to earlier outputs. Projections have depth zero;
finite unions take the maximum depth of their arguments. On negatives
using each color at most once, the least depth needed for a positive
query has the following description.

\begin{theorem}[Minimum completion depth]\label{core:thm:completion-depth}
Suppose $|t_r|\le1$ on every negative. Starting at $S=Z$, the minimum
depth for $(i,p)$ is the minimum number of moves needed to reach
$i\notin h(S)$, with minimum $\infty$ when this is impossible.
The moves are:
\begin{enumerate}
\item choose an owner $r$ and delete all rows with
      $t_r\setminus D_r(S;p)\ne\varnothing$;
\item choose $v\notin h(S)$ and delete all rows posting $v$ at a
      color $r$ for which $v\notin p_r$.
\end{enumerate}
\end{theorem}
\begin{proof}
At a terminal budget take $x_i$. For the first move, combine a
witness on the remaining rows and all inputs outside $D_r(S;p)$ in
one $r$-fold. Such a fold is its greatest completion: it absorbs an
$r$ singleton and otherwise takes the union. The reachability proof
of \cref{core:thm:injective-blockers} shows that it fits the budget
and passes the mark. For the second move use a binary completion on
the witness and $x_v$. Give it the second projection's profile at the
colors $r$ with $v\notin p_r$, and unanimity at every other color.
Every winning coalition contains the second input, so these profiles
are admissible. The completion absorbs on the deleted rows.
Elsewhere $x_v$ supplies no negative mark, and at the positive no
profile absorbs: the witness is ambiguous and $v$ is unposted at
each projection color.

Conversely take a fitting expression. A free union passes to a
mark-supplying summand that still omits $S$. At a completion root,
each row with singleton output $\{r\}$ has $t_r=\{v\}$ and
$\Phi_r(\{v\})=1$ for the whole expression. All such $v$ lie outside
$D_r(S;p)$. Indeed the mark precludes a positive bit at $p_r$.
A positive bit on a negative head forces the $r$ singleton there
by \cref{core:lem:singleton-profile-absorption}, as the alternative
contains its head menu. Thus empty $r$-tails have losing heads and a
losing singleton tail has a losing head, proving the claim by reachability.
If all singleton outputs have one color, the first move covers them.
If two or more colors occur, every winning singleton coalition just
identified must use the same input $v$, by the rainbow intersection
condition. Two different colors winning on $\{v\}$ force the entire
menu of $v$ to survive at every valuation, again by singleton-profile
absorption. Hence $v\notin h(S)$, and the second move covers all
these rows. Wherever the completion is ambiguous it takes the union
of its supported children. One such child supplies the positive, and
each fits the remaining rows. If there were no singleton outputs one
can pass to that child without a move. Descent gives a schedule of
length at most the expression's depth.
\end{proof}

\begin{corollary}[Sharp depth and shortest common supersequences]
\label{core:cor:completion-depth-hierarchy}
For every $m\ge1$ the maximum least depth of a feasible sample with
$m$ negatives and one positive is $m$. The classes of operations of
successive depths are strictly increasing for every $K\ge3$.
On explicitly listed negatives with $|t_r|\le1$, deciding whether
depth at most $d$ suffices is NP-complete at every fixed $K\ge3$.
\end{corollary}
\begin{proof}
Even for unrestricted negatives the two-child construction of
\cref{core:thm:pointed-guard-cut} gives depth at most $m$: each owner
choice is a binary greatest completion and drops a row on both branches.
For the lower bound take a path $x_0,\ldots,x_m$. Its $j$th edge gives
a negative headed at $x_{j-1}$ with sole posting $x_j:r_j$; ask for
$x_0$ at the empty coloring. More generally, take several such paths
sharing only $x_0$. In a remaining path prefix an owner move deletes
its entire terminal run of that color, and no earlier edge: the
terminal target is head-free, while the first different color puts
its head in the initial set for reachability, which then propagates
backwards through the preceding run. A head-free input move deletes
at most one terminal edge, since targets are private to paths.
Replacing it with the corresponding owner move can only shorten the
remaining prefixes. It therefore suffices to minimize over owner
schedules. Such a schedule succeeds exactly when it contains every
reversed run sequence as a subsequence. The least depth is the
length of a shortest common supersequence.

One path of alternating colors gives depth $m$. Its left-associated
word obtained by adjoining $x_j$ by $\star_{r_j}$ at step $j$ fits
this sample, so its full operation also has exact depth $m$.
For the complexity claim, a schedule of at most $\min(d,|Z|)$
productive moves is a polynomially checkable certificate. Conversely
use one path for the reverse of each input word of the three-letter
shortest common supersequence problem with no consecutive equal
letters. That problem is NP-complete
\cite[Corollary~6]{LagoutteTavenas2015}, and the construction preserves
the optimum. Additional owners delete no rows.
\end{proof}

\begin{theorem}[Sharp arity bound for completion depth]
\label{core:thm:additive-completion-depth}
For $n\ge1$ and $K\ge3$, the greatest depth required by a feasible
sample on $n$ inputs with one positive, when each color occurs at most
once in each negative, is
\[
                         \left\lfloor\frac{3(n-1)}2\right\rfloor .
\]
The same maximum holds for realizable complete families under this
restriction on their minimal losses. It is attained even by families
with one prime guard per observer and source and one posting per loss.
For $n\ge2$, restricting all postings to two colors lowers the maximum
to $n-1$; restricting them to one color lowers it to one.
\end{theorem}
\begin{proof}
Use the two moves of \cref{core:thm:completion-depth}, stopping as
soon as the queried input $i$ is head-free. Whenever a head-free
input move can remove a row, take such a move. Otherwise take a
productive owner move. These choices preserve feasibility since
they only delete rows. A nonterminal feasible budget always admits
a productive move by the schedule theorem.

A head-free input move at $v$ clears every posting of $v$ eligible
at the positive, so it can be used only once. Now consider an owner
move at $r$, and choose a removed posting $v\notin D_r(S;p)$.
Since $v\notin p_r$, the priority rule implies that $v$ is still a
head. Every row headed at $v$ must have a posting $u$ at $r$ outside
$D_r(S;p)$: an empty tail would seed $v$, and a reached tail would
reach it. Choose such a row. Its posting $u$ differs from its head
$v$, and the same priority argument makes $u$ a head. Every row
headed at either $u$ or $v$ has its $r$-posting outside $D_r(S;p)$,
so this move frees at least two previously occupied heads.

If the last move is an input move, no owner move freed $i$. There
are at most $\lfloor(n-1)/2\rfloor$ owner moves and at most $n-1$
input moves. If the last move is an owner move, its two newly freed
heads are never used for input moves. There are then at most
$\lfloor n/2\rfloor$ owner moves and at most $n-2$ input moves.
Both totals give the upper bound; an initially head-free $i$ needs
no move. Taking the free union of witnesses over the positive bits
gives the same bound for complete families.

With only two posting colors, continue to take productive head-free
moves first. Each such move settles one input: it has no eligible
posting left. At an owner move at $r$ choose the two newly freed
inputs $u,v$ above. All their $r$-postings have been removed. If $i$
is now free, charge just this step to one of $u,v$ different from $i$
and stop. Otherwise clear the eligible postings of $u,v$ at the other
color $s$ in one further layer, if needed. Attach an $s$-fold of the
residual witness and those of $u,v$ outside $p_s$ that still post at
$s$. It absorbs on the deleted rows, misses the mark at the heads of
all retained rows since its added inputs are head-free, and preserves
the positive. This settles both inputs. No settled input can be charged
again, so at most $n-1$ layers are charged before $i$ is free.
With only one posting color $r$, the owner move clears all rows with
a tail outside $D_r(S;p)$. Its deletion preserves $D_r(S;p)$: the
initial heads and every edge between reached vertices remain. No
further owner move can then remove a row; an eligible head-free tail
would itself be outside the reached set. By feasibility the queried
head must already be free.

For sharpness take $d=\lfloor(n-1)/2\rfloor$ ordered pairs
$\{u_j,v_j\}$ followed by a final input $i$. For even $n$ put one
more singleton $h$ immediately before $i$. In each pair put both
internal arrows at $a$. To every later vertex put $a$-arrows from
both members, $b$-arrows from $u_j$ and $c$-arrows from $v_j$, with
$a,b,c$ distinct. Add $h\to_a i$ if $h$ is present. With reflexive
arrows understood these data satisfy descent and reciprocity: every
nonloop first arrow already reaches every proper descendant in its
first color, and reverse arrows occur only within a pair.
The reconstruction is core. On each observer facet its sole prime
excludes the incoming postings visible there. Each such literal
is necessary on its one-posting row.

Ask for $i$ at the empty coloring. At the first remaining pair the
only productive initial move is at owner $a$: its two internal heads
seed the other colors, and its foreign arrows seed every later
vertex at $a$. That move clears the pair's $a$-arrows. It leaves a
$b$-posting from $u_j$ and a $c$-posting from $v_j$ to every later
vertex, including $i$. Clearing one kind requires its owner move
or the head-free move at its own member. Neither clears the other
kind. Until both are cleared, every later pair is still seeded at
$a$; its internal $a$-arrows seed $b,c$ and keep both members from
being head-free. No later pair can begin early. Thus each pair costs
three moves in succession, and $h\to_a i$ costs a further move in
the even case. This gives $3d+[n\text{ even}]$ as claimed.

For two colors take one directed path of $n$ vertices with alternating
outgoing colors. Give each vertex arrows of its outgoing color to all
later vertices. This is a realizable conjunctive family by descent
and reciprocity. Its adjacent-edge sample alone requires $n-1$ moves
by \cref{core:cor:completion-depth-hierarchy}. For one color a single
nonloop arrow already requires one layer.
\end{proof}

\begin{corollary}[Optimizing a complete conjunctive guard family]
\label{core:cor:additive-depth-complexity}
For a realizable complete family with one prime per observer and
source, every positive coloring at a given source has the same least
witness depth as its empty coloring. The depth is computable in
polynomial time with at most two posting colors. With three posting
colors, deciding depth at most $d$ is NP-complete for each fixed
$K\ge3$, already when the capture graph has no nonloop cycle.
\end{corollary}
\begin{proof}
Fix source $i$ and its predecessor set $H$, including $i$.
Restricting a deletion schedule to rows headed in $H$ cannot increase
its length. Reachability can only shrink, and an input move outside
$H$ clears no such row. Conversely a witness on $H$ ignores every
other head. For every nonloop arrow $u\to_r v$ within $H$, either
$u=i$ or descent along a path from $v$ to $i$ gives $u\to_r i$.
Thus a positive $p$ posts no $r$-tail of any such row at $r$.
In every remaining budget the extra seeds $p_r$ in $D_r(S;p)$
therefore have no outgoing reachability edges and are not candidate
$r$-tails. They change neither the rows removed by an owner move nor
the permitted postings of a head-free input move. The possible
schedules on $H$ are exactly those for the empty coloring.

The depth of the operation is the maximum of these minima over the
heads: the operation witnesses each query, and the free union of
witnesses of all its positive bits realizes it. One schedule per head
is therefore an NP certificate, polynomial in size by
\cref{core:thm:additive-completion-depth}.

With two colors $b,c$, an owner move at $r$ leaves $D_r$ unchanged:
it deletes only edges from outside $D_r$, and no empty $r$-tail.
After this move every currently head-free input has no remaining
$r$-posting. Indeed it lies outside $D_r$ on the residual graph,
as $p$ is empty. Any subsequent input move can therefore be replaced
by the other owner move, which deletes at least the same rows.
All moves remain valid on smaller budgets, with a possibly larger
set of deletions. Repeated moves of the same owner are redundant.
Thus one may alternate owners after the first owner move.

Among optimal schedules of this form minimize the length of the
initial input-only prefix. Two consecutive inputs $v,w$ in the
prefix cannot both already be free before $v$: an owner move at
$b$ followed by one at $c$ covers all postings of both, permitting
an earlier switch. Thus $v$ frees $w$, which requires an arrow
$v\to w$. The chosen inputs form a chain in capture order. At this
stage only postings out of already chosen inputs have been deleted,
so every proper predecessor of a chosen input was chosen earlier.
A nontrivial component cannot supply the first of its own inputs.
Consequently a nonempty prefix is exactly the predecessor set of
its final vertex, including that vertex, linearly ordered with
singleton components. Enumerate these at most $|H|$ prefixes and the
empty one. After each try both alternating owner sequences until
$i$ is free or a complete pair of moves makes no progress. Taking
the least length over these polynomially many simulations gives
the exact depth.

For hardness take the disjoint paths ending at a common queried head
from \cref{core:cor:completion-depth-hierarchy}, oriented from a
posting to its head. On each path, give every vertex of outgoing color
$r$ an $r$-arrow to every vertex strictly closer to the common head.
There are quadratically many arrows. Descent and reciprocity hold,
so their conjunctive guards give a realizable operation. Each
head-free input move deletes postings of just one color and is
dominated by that owner move. On a path an owner move clears its
initial run of the chosen color in the direction of the arrows.
A vertex beyond an earlier remaining vertex of another color is
already seeded for this move by the direct arrow between them.
Thus it cannot delete postings out of order, even after completion
of the paths. The least depth at the common head is exactly the
shortest common supersequence length. Other heads use just a prefix
of one path in that order and require no greater depth. The optimum
of the complete operation is therefore the same as that of the
original supersequence instance.
\end{proof}

\begin{proposition}[Exponential length of fixed offer lists]
\label{core:prop:fixed-offer-parity}
At each fixed $K\ge3$ there are realizable complete prime families on
$O(d)$ inputs, of polynomial description size with at most four postings
per guard, and a positive query admitting a winning list of arbitrary
coalition offers fixed before any responses. Every such list has length
at least $2^{d-1}-1$.
\end{proposition}
\begin{proof}
Use the directed cube masks of \cref{core:prop:compact-cube-masks}.
A bit string $x$ selects one rail per bit, forming $L_x$. An edge
$e:x\to y$ masks $(L_x,U_e)$ with $U_e=L_x\cup L_y$.
Use its opposite probe: its head is the rail in $L_x\setminus L_y$,
with $p_c=L_y$ and $p_b=V\setminus U_e$, observed by a third owner.
If only one direction of each edge occurs, this probe is positive.

By the forward proof of \cref{core:cor:fixed-offer-prime-space}, a
winning list of length $s$ gives a monotone Boolean DNF $f$ with at
most $s+1$ terms. Pull it back from $H$ to all the rails by intersection
with $H$. It satisfies $f(L_y)=0$, $f(U_e)=1$, and is constant on
every masked cover whose differing rail belongs to $H$. On the
other covers the two endpoints have the same intersection with $H$.
Consequently $f(L_v)\le f(L_u)$ on every directed edge $u\to v$.

Now store an address $x\in\{0,1\}^d$. Interpret its first $d-1$ bits
as the low bits of a binary counter and the parity of all $d$ bits as
its high bit. Increment in these coordinates, stopping at the last
count. Encoding, binary increment and decoding use $O(d)$ two-input
gates. Starting at zero, the counter visits every even-parity address
before every odd one. Apply the one-bit circuit sweep from
\cref{core:prop:compact-cube-masks}, testing the continue bit before
entering a sequence of constant-zero writes that clears all scratch
before returning to the first stage. Each scratch bit now has at most
two writing stages, one subcube per direction at each. Each directed
control flip still has at most two subcubes. The same prime calculation
therefore gives the stated description bound. The directed graph is
acyclic: a cycle would pass the first stage and hence would repeatedly
advance a counter that halts. Its complete mask is core.

Choose $e$ on the run between the last even and first odd address at
the first stage. There the scratch bits are zero and the flags fixed.
The gap above gives $f(L_x)=1$, $f(L_y)=0$ at the ends of $e$, so on
all the first-stage addresses $f$ is precisely even parity.
Each rail term restricts there to a conjunction of address literals.
A satisfiable conjunction implying parity fixes every bit, so at least
$2^{d-1}$ terms are necessary.

A fixed winning list exists at this positive by the acyclic case in
\cref{core:cor:fixed-offer-prime-space}.
\end{proof}

\paragraph{Nested repeated coalitions of one owner.}
Suppose all pairs $(r,t_r)$ with $|t_r|\ge2$ have the same owner $a$
and that these coalitions form a chain under inclusion. Arbitrarily
many negative rows may use any of them, mixed with other colors.
Fix a positive query $(i,p)$ and put $B=p_a$. Use the marked test $M$
and the profile-gap test $G(A)$ asking for $\Phi_a(A)=1$ and
$\Phi_a(B)=0$. For either test on rows $S$, call $S$ \emph{closed
for the test} when every eligible projection is a head of $S$ and
$t_r\subseteq D_r$ for every row and owner, with the seeds of that
test in \cref{core:thm:repeated-blocker-budgets}. The eligible
projections are $i$ for $M$ and the vertices of $A\setminus B$
for $G(A)$.

\begin{theorem}[A flag criterion for nested repetitions]
\label{core:thm:nested-repeated-guards}
The positive query is impossible exactly when there are a flag
$S_0\supsetneq\cdots\supsetneq S_h$ of subfamilies of the given
negatives and a strictly decreasing chain $Q_0\supsetneq\cdots
\supsetneq Q_{h-1}$ of repeated coalitions with the following properties.
\begin{enumerate}
\item $S_0$ is closed for $M$, and $S_j$ is closed for $G(Q_{j-1})$
for $1\le j\le h$.
\item At each $j<h$, the largest repeated coalition in $S_j$ is $Q_j$,
and no row of $S_{j+1}$ contains $Q_j$ in its $a$-class.
The last family $S_h$ has no repetition.
\end{enumerate}
The test takes polynomial time on explicitly listed queries. With $n$
inputs a failure already occurs on at most $n(n\min(K,n)+1)$ negative
queries.
\end{theorem}
\begin{proof}
At a gap test discard all rows containing $A$ in their $a$-class:
a witness for the gap absorbs on all of them. At either test apply
the projection and fold rules until they succeed or leave a closed
subfamily. At a closed subfamily without a repetition the test fails.
Otherwise let $Q$ be its largest repetition. By
\cref{core:thm:repeated-blocker-budgets}, feasibility requires a
witness winning at some repeated coalition and hence, by inclusion,
one winning at $Q$. Put $S'=\{(j,t)\in S:Q\not\subseteq t_a\}$.

For $M$, feasibility is equivalent to both $M$ and $G(Q)$ being
feasible on $S'$. A marked witness winning at $Q$ also
loses at $B$ because it supplies the mark. Conversely attach a gap
witness to a marked witness by $\star_a$. Both omit the marks on
$S'$, the gap witness absorbs on every removed row, and its loss at
$B$ protects the positive mark. For $G(A)$ with $A$ in the chain,
the initial deletion gives $Q\subsetneq A$. Winning at $Q$ supplies
the old demand, so it suffices to test $G(Q)$ alone on $S'$.

If the mark is impossible, take an inclusion-minimal failing subfamily
$S_0$. It must be closed, since a fold would give a smaller failing
subfamily. If it has a repeat, $M$ is feasible on its proper subset
$S'$, so $G(Q_0)$ fails there. Successively take a minimal failing
subfamily for this gap and the resulting smaller gaps. This gives
the flag. Conversely its terminal family fails. At every preceding
family the next gap is necessary already on the full set $S'$.
Failure on the displayed subfamily therefore implies failure one
stage earlier.

For polynomial time one need not find minimal subfamilies. Apply
the projection, fold and largest-coalition tests just proved. Each
successive largest coalition is strictly smaller. There are at most
$n$ marked stages, each starting at most one sequence of at most $n$
gap stages. Projection and reachability take polynomial time.
To shorten a failing flag, retain at each stage one row per head.
If its $r$-tail is a singleton and that head is outside $L_r$, also
retain the row first reaching the head in the original $r$-search;
otherwise the head is already a seed or belongs to $L_r$.
As in \cref{core:cor:input-guard-cuts}, at most $n\min(K,n)$ rows preserve
all heads and reachable sets. Retain one further row for the largest
repetition if present. At each stage also retain the witnesses from
all later stages. They were already in that stage, so its closure
and largest repetition persist. There are at most $n$ stages.
\end{proof}
